\documentclass[11pt,a4paper]{article}

% ---------- Codificación y tipografía ----------
\usepackage[utf8]{inputenc}
\usepackage[T1]{fontenc}
\usepackage{lmodern}
\usepackage[spanish,es-nodecimaldot]{babel}
\usepackage{microtype}
\usepackage{amsmath,amssymb}
\usepackage{geometry}
\geometry{top=2.3cm,bottom=2.3cm,left=2.2cm,right=2.2cm}
\usepackage{enumitem}
\usepackage{booktabs}
\usepackage{longtable}
\usepackage{array}
\usepackage{fancyhdr}
\usepackage{titlesec}
\usepackage{setspace}
\usepackage{multicol}
\usepackage{needspace}

\setstretch{1.08}
\setlength{\parindent}{0pt}
\setlength{\parskip}{4pt}

% ---------- Encabezado y pie ----------
\pagestyle{fancy}
\fancyhf{}
\lhead{\small Banco de Preguntas PEU --- SUNASS}
\rhead{\small Módulo IX: Econometría}
\cfoot{\small \thepage}
\renewcommand{\headrulewidth}{0.4pt}
\renewcommand{\footrulewidth}{0pt}

% ---------- Títulos ----------
\titleformat{\section}{\large\bfseries}{\thesection}{0.8em}{}
\titlespacing*{\section}{0pt}{14pt}{6pt}
\titleformat{\subsection}{\normalsize\bfseries}{\thesubsection}{0.6em}{}
\titlespacing*{\subsection}{0pt}{10pt}{4pt}

% ---------- Comandos del examen ----------
\newcounter{qnum}
\newcommand{\q}[1]{%
  \needspace{8\baselineskip}%
  \par\medskip\refstepcounter{qnum}%
  \noindent\textbf{\theqnum.}\ #1\par\nopagebreak[4]
}
\newenvironment{opts}
  {\begin{enumerate}[label=\Alph*),leftmargin=2.0em,itemsep=1pt,topsep=2pt,parsep=0pt,partopsep=0pt]}
  {\end{enumerate}}
\makeatletter
\@itempenalty=10000
\makeatother
\newcommand{\sol}[3]{\par\noindent\textbf{#1. #2.}\ #3\par\smallskip}

% ==============================================================================
% Configuración BackBone Academy para los Bancos de Preguntas (Módulos 01-10)
% ==============================================================================
\usepackage{graphicx}
\usepackage{xcolor}
\usepackage{tikz}
\usepackage[hidelinks]{hyperref}

% ==============================================================================
% Identidad de marca BackBone Academy compartida (Beamer y Bancos de Preguntas)
% ==============================================================================
\definecolor{azulPrincipal}{RGB}{43, 45, 162}
\definecolor{azulOscuro}{RGB}{24, 26, 98}
\definecolor{grisTexto}{RGB}{85, 85, 95}

\newcommand{\backboneIgUrl}{https://www.instagram.com/backbone.academy}
\newcommand{\backboneTtUrl}{https://www.tiktok.com/@backbone.academy}
\newcommand{\backboneWebUrl}{https://leonardovelasquez14-detg.github.io/peu-sunass-economia/}

% Icono discreto de Instagram en TikZ
\newcommand{\igicon}[1][white!88]{%
  \tikz[baseline=-0.25ex]{%
    \draw[#1, line width=0.65pt, rounded corners=1.1pt] (0,0) rectangle (0.23cm,0.23cm);
    \draw[#1, line width=0.6pt] (0.115cm,0.115cm) circle (0.055cm);
    \fill[#1] (0.18cm,0.18cm) circle (0.014cm);
  }%
}

% Icono discreto de TikTok en TikZ (nota musical estilizada)
\newcommand{\tiktokicon}[1][white!88]{%
  \tikz[baseline=-0.25ex]{%
    \draw[#1, line width=0.82pt, line cap=round, line join=round]
      (0.11cm,0.06cm) -- (0.11cm,0.22cm)
      .. controls (0.12cm,0.16cm) and (0.16cm,0.145cm) .. (0.20cm,0.145cm);
    \draw[#1, line width=0.82pt, line cap=round]
      (0.11cm,0.065cm) arc[start angle=0, end angle=295, radius=0.048cm];
  }%
}

% Enlaces compactos de redes sociales: \backboneSocials{color}{tamaño}
\newcommand{\backboneSocials}[2]{%
  #2\color{#1}%
  \href{\backboneIgUrl}{\igicon[#1]}%
  \hspace{3.5pt}%
  \href{\backboneTtUrl}{\tiktokicon[#1]}%
  \hspace{4pt}%
  \href{\backboneIgUrl}{\textbf{@backbone.academy}}%
}


% Pie de página institucional en todas las páginas interiores
\lfoot{\scriptsize\color{azulOscuro}\textbf{BackBone Academy}\ \color{grisTexto}· \href{\backboneIgUrl}{\igicon[azulPrincipal]\ \textbf{@backbone.academy}} · \href{\backboneTtUrl}{\tiktokicon[azulPrincipal]\ \textbf{TikTok}} · \href{\backboneWebUrl}{Simulador Web}}
\cfoot{}
\rfoot{\small\thepage}
\renewcommand{\footrulewidth}{0.3pt}

% Encabezado superior para la portada de cada banco
\newcommand{\backboneportadaheader}{%
  \vspace*{0.1cm}
  \noindent\includegraphics[height=0.72cm]{../teoria/backbone_logo_azul.png}%
  \hfill%
  \raisebox{0.22cm}{\backboneSocials{azulPrincipal}{\small}}%
  \par\vspace{0.9cm}%
}

% Pie institucional para la portada de cada banco
\newcommand{\backboneportadafooter}{%
  {\small\color{azulOscuro}\textbf{BackBone Academy} \color{black}· Material de preparación académica.}\\[0.18cm]
  {\small\color{azulPrincipal}\href{\backboneIgUrl}{\igicon[azulPrincipal]\ \textbf{Instagram}} \quad·\quad \href{\backboneTtUrl}{\tiktokicon[azulPrincipal]\ \textbf{TikTok}} \textbf{(@backbone.academy)} \quad·\quad \href{\backboneWebUrl}{\textbf{Simulador Web Interactivo}}}\\[0.18cm]
  {\small Edición 2026}%
}


\begin{document}

% ============================================================
% PORTADA
% ============================================================
\begin{titlepage}
\thispagestyle{empty}
\centering
\backboneportadaheader

{\Large\bfseries BANCO DE PREGUNTAS}\\[0.35cm]
{\large EXAMEN DE ADMISIÓN}\\[0.15cm]
{\large PROGRAMA DE EXTENSIÓN UNIVERSITARIA (PEU)}\\[0.15cm]
{\large SUPERINTENDENCIA NACIONAL DE SERVICIOS DE SANEAMIENTO}\\[0.05cm]
{\large SUNASS}\\[1.0cm]

\rule{\linewidth}{0.9pt}\\[0.55cm]
{\LARGE\bfseries MÓDULO IX: ECONOMETRÍA}\\[0.35cm]
{\large Especialidad: Economía}\\[0.55cm]
\rule{\linewidth}{0.9pt}\\[1.0cm]

\begin{tabular}{ll}
\toprule
Número de preguntas: & 180 \\
Bloque A: & 30 preguntas \\
Bloque B: & 35 preguntas \\
Bloque C: & 25 preguntas \\
Bloque D: & 90 preguntas (63 teóricas, 27 prácticas) \\
Duración sugerida: & 150 minutos \\
Nivel: & Econometría intermedia--avanzada \\
Formato: & Opción múltiple (A, B, C, D) \\
\bottomrule
\end{tabular}

\vfill
\backboneportadafooter
\end{titlepage}

% ============================================================
% INSTRUCCIONES
% ============================================================
\section*{Instrucciones generales}

\begin{enumerate}[label=\arabic*.,leftmargin=1.6em,itemsep=3pt]
  \item El presente banco contiene \textbf{180 preguntas de opción múltiple}. Cada pregunta tiene una sola respuesta correcta.
  \item No se descuenta por respuesta incorrecta; se recomienda responder todas las preguntas.
  \item El tiempo sugerido de resolución es de \textbf{150 minutos}. Use calculadora no programable.
  \item Las preguntas están organizadas en cuatro bloques: fundamentos teóricos, aplicaciones cuantitativas y casos regulatorios del sector saneamiento.
  \item El \textbf{solucionario completo} y la \textbf{tabla de claves rápidas} se encuentran al final. Se recomienda resolver el examen antes de consultarlos.
\end{enumerate}

\section*{Estructura del banco}

\begin{center}
\renewcommand{\arraystretch}{1.25}
\begin{tabular}{@{}p{1.4cm}>{\raggedright\arraybackslash}p{7.6cm}cc@{}}
\toprule
\textbf{Bloque} & \textbf{Contenido} & \textbf{Preguntas} & \textbf{Peso} \\
\midrule
A & Fundamentos teóricos: modelo clásico de regresión lineal, estimador MCO, supuestos de Gauss-Markov, Teorema BLUE, inferencia ($t, F$), $R^2$, multicolinealidad, heterocedasticidad, autocorrelación y endogeneidad. & 1--30 & 16.7\% \\
B & Aplicaciones cuantitativas: cálculo de estimadores, estadísticos $t$ y $F$, contrastes de White, Breusch-Pagan y Durbin-Watson, MCG/WLS, sesgo de variable omitida, variables instrumentales (MC2E / 2SLS) y diagnósticos. & 31--65 & 19.4\% \\
C & Casos con implicancias regulatorias: estimación de demanda de agua y elasticidades, tarifas por bloques y endogeneidad, fronteras estocásticas de costos (SFA) en EPS, diferencias en diferencias (Diff-in-Diff) y WACC regulatorio. & 66--90 & 13.9\% \\
D & Ampliación integral del módulo: 63 preguntas teóricas y conceptuales (70\%) y 27 aplicaciones prácticas cuantitativas (30\%). & 91--180 & 50.0\% \\
\bottomrule
\end{tabular}
\end{center}

\section*{Competencias evaluadas}

\noindent El banco evalúa la capacidad de comprender los fundamentos estadísticos y matemáticos del Modelo Clásico de Regresión Lineal (MCO), verificar el cumplimiento de los supuestos de Gauss-Markov, realizar inferencia estadística robusta mediante pruebas de hipótesis individuales y conjuntas, diagnosticar y corregir violaciones clásicas (multicolinealidad, heterocedasticidad, correlación serial y endogeneidad) y aplicar métodos econométricos rigurosos a la modelación de la demanda de agua potable, funciones de costos de EPS, evaluación de impacto y determinación del costo de capital regulatorio.

\section*{Notación principal}

{\small
\renewcommand{\arraystretch}{1.0}
\begin{center}
\begin{tabular}{@{}ll@{\hspace{1.0cm}}ll@{}}
\toprule
$\hat{\beta}$ & Estimador de Mínimos Cuadrados Ordinarios & $X$ & Matriz de regresores ($n \times k$) \\
$u, e$ & Perturbación estocástica y residuo muestral & $\sigma^2$ & Varianza del término de error \\
$R^2, \bar{R}^2$ & Coeficiente de determinación simple y ajustado & $VIF$ & Factor de inflación de la varianza \\
$t, F$ & Estadísticos de prueba t-Student y F-Snedecor & $d$ & Estadístico de Durbin-Watson \\
$MCG / WLS$ & Mínimos cuadrados generalizados / ponderados & $MC2E / 2SLS$ & Mínimos cuadrados en dos etapas \\
$HAC$ & Errores estándar robustos de Newey-West & $SFA$ & Análisis de frontera estocástica \\
\bottomrule
\end{tabular}
\end{center}
}

\newpage
\setcounter{qnum}{0}

\section*{Bloque A. Fundamentos teóricos y conceptuales}

\noindent\textit{Preguntas 1 a 30. Evalúan la comprensión rigurosa del modelo clásico de regresión lineal, estimador MCO, supuestos de Gauss-Markov, Teorema BLUE, inferencia estadística ($t, F$), $R^2$, modelos no lineales, multicolinealidad, heterocedasticidad, autocorrelación y endogeneidad.}

\q{En la representación matricial del Modelo Lineal General $Y = X\beta + u$, el estimador de Mínimos Cuadrados Ordinarios (MCO) $\hat{\beta}$ se obtiene minimizando la suma de residuos al cuadrado y su expresión cerrada es:}
\begin{opts}
    \item $\hat{\beta} = (X'X) X'Y$
    \item $\hat{\beta} = (X'X)^{-1} X'Y$
    \item $\hat{\beta} = (X'Y)^{-1} X'X$
    \item $\hat{\beta} = X'(XX')^{-1} Y$
\end{opts}

\q{El supuesto de 'Exogeneidad Estricta' en el modelo de regresión lineal establece formalmente que $E[u|X] = 0$. La consecuencia teórica directa más importante de este supuesto es que:}
\begin{opts}
    \item el estimador MCO es asintóticamente eficiente pero presenta sesgo en muestras finitas.
    \item los errores estándar de los estimadores convergen a cero cuando el tamaño de muestra es finito.
    \item el coeficiente de determinación $R^2$ alcanza siempre un valor unitario en muestras finitas.
    \item el estimador MCO es insesgado respecto al vector de parámetros poblacionales ($E[\hat{\beta}|X] = \beta$).
\end{opts}

\q{Para que la matriz $(X'X)^{-1}$ exista y el estimador MCO pueda calcularse unívocamente, se requiere el supuesto de:}
\begin{opts}
    \item normalidad estricta de la variable dependiente $Y$ en las observaciones.
    \item homocedasticidad estricta y ausencia de autocorrelación serial en los errores.
    \item rango de columna completo de la matriz $X$, descartando multicolinealidad perfecta.
    \item independencia temporal completa entre las perturbaciones estocásticas.
\end{opts}

\q{El supuesto de perturbaciones esféricas en el modelo clásico de Gauss-Markov abarca dos condiciones fundamentales sobre la matriz de varianzas-covarianzas condicional del error $Var(u|X) = \sigma^2 I_n$:}
\begin{opts}
    \item estacionariedad en sentido estricto y ergodicidad asintótica del término de error.
    \item normalidad bivariada, asimetría nula y curtosis exacta igual a tres en los residuos.
    \item homocedasticidad (varianza constante $\sigma^2$) y ausencia de autocorrelación serial.
    \item independencia estocástica lineal entre los regresores y la variable dependiente medida.
\end{opts}

\q{El célebre Teorema de Gauss-Markov demuestra que, bajo los supuestos clásicos de regresión lineal, el estimador MCO:}
\begin{opts}
    \item coincide con el estimador de máxima verosimilitud siempre que los errores sigan distribución uniforme.
    \item alcanza un coeficiente $R^2$ superior al de cualquier estimador no lineal bajo errores esféricos.
    \item es invariante en magnitud frente a cualquier transformación no lineal en la escala de las variables.
    \item es el Mejor Estimador Lineal Insesgado (MELI), con varianza mínima entre estimadores lineales insesgados.
\end{opts}

\q{El estimador insesgado de la varianza poblacional de las perturbaciones estocásticas ($\sigma^2$) en una regresión múltiple con $n$ observaciones y $k$ parámetros (incluyendo la constante) es:}
\begin{opts}
    \item $s^2 = \frac{\sum_{i=1}^n e_i^2}{n - 1}$
    \item $s^2 = \frac{\sum_{i=1}^n e_i}{n - k}$
    \item $s^2 = \frac{e'e}{n - k}$
    \item $s^2 = \frac{\sum_{i=1}^n e_i^2}{n}$
\end{opts}

\q{El coeficiente de determinación $R^2$ mide la proporción de la variación total de la variable dependiente explicada por el modelo de regresión. Una propiedad matemática intrínseca del $R^2$ es que:}
\begin{opts}
    \item disminuye necesariamente al introducir cualquier regresor adicional irrelevante.
    \item su valor siempre oscila entre $-1$ y $+1$ cuando el modelo incluye constante.
    \item nunca disminuye al agregar variables explicativas adicionales, aunque no sean significativas.
    \item es igual a cero solo cuando todos los coeficientes estimados son negativos.
\end{opts}

\q{Para penalizar la inclusión innecesaria de regresores que no aportan poder explicativo genuino, se utiliza el coeficiente $R^2$ ajustado ($\bar{R}^2$), el cual se define como:}
\begin{opts}
    \item $\bar{R}^2 = 1 - \frac{SCR}{STC} \sqrt{n - k}$
    \item $\bar{R}^2 = 1 - \frac{n - k}{n - 1} (1 - R^2)$
    \item $\bar{R}^2 = 1 - \frac{SEC / (n - k)}{STC / (n - 1)}$
    \item $\bar{R}^2 = 1 - \frac{SCR / (n - k)}{STC / (n - 1)}$
\end{opts}

\q{Si además de los supuestos de Gauss-Markov se añade el supuesto de Normalidad de los errores ($u|X \sim N(0, \sigma^2 I_n)$), la distribución muestral exacta del estimador MCO es:}
\begin{opts}
    \item $\hat{\beta}|X \sim N(\beta, \sigma^2 (X'X)^{-1})$
    \item $\hat{\beta}|X \sim t_{n-k}(\beta, \sigma^2 I_n)$
    \item $\hat{\beta}|X \sim F_{k, n-k}(0, \sigma^2 I_n)$
    \item $\hat{\beta}|X \sim \chi^2_{n-k}(\beta, \sigma^2 X'X)$
\end{opts}

\q{Para contrastar la hipótesis nula de significancia individual $H_0: \beta_j = 0$ contra $H_1: \beta_j \neq 0$, se utiliza el estadístico $t$ de Student definido como:}
\begin{opts}
    \item $t = \frac{SCR}{SEC}$, que bajo $H_0$ sigue una distribución normal estándar en muestras finitas.
    \item $t = \frac{\hat{\beta}_j - 1}{EE(\hat{\beta}_j)}$, con distribución normal estándar para cualquier tamaño de muestra.
    \item $t = \frac{\hat{\beta}_j}{EE(\hat{\beta}_j)}$, que bajo $H_0$ sigue una $t$ con $n - k$ grados de libertad.
    \item $t = \frac{\hat{\beta}_j}{s^2}$, que bajo $H_0$ sigue una distribución $\chi^2$ con $n - k$ grados de libertad.
\end{opts}

\q{La prueba $F$ de significancia conjunta del modelo contrasta la hipótesis nula $H_0: \beta_2 = \beta_3 = \dots = \beta_k = 0$ (todos los coeficientes de las pendientes son cero simultáneamente) frente a la alternativa de que al menos uno es distinto de cero. Su fórmula a partir del $R^2$ es:}
\begin{opts}
    \item $F = t^2 \times \frac{n - k}{k - 1}$, con distribución $t_{n-k}$ bajo $H_0$.
    \item $F = \frac{1 - R^2}{R^2} (n - k)$, con distribución $F_{n-k, k-1}$ bajo $H_0$.
    \item $F = \frac{R^2 / (k - 1)}{(1 - R^2) / (n - k)}$, con distribución $F_{k-1, n-k}$ bajo $H_0$.
    \item $F = \frac{R^2 / (n - k)}{(1 - R^2) / (k - 1)}$, con distribución $\chi^2_{k-1}$ bajo $H_0$.
\end{opts}

\q{En un modelo econométrico con especificación doble logarítmica (log-log) $\ln(Y_i) = \beta_0 + \beta_1 \ln(X_i) + u_i$, el coeficiente de la pendiente $\beta_1$ se interpreta directamente como:}
\begin{opts}
    \item la tasa de crecimiento porcentual instantánea de la constante ante variaciones unitarias en $X$.
    \item la elasticidad constante de $Y$ respecto a $X$ (variación porcentual de $Y$ por $1\%$ en $X$).
    \item la variación absoluta en niveles de $Y$ ante una variación unitaria en el logaritmo de $X$.
    \item la semielasticidad de $Y$ ante un cambio discreto de una unidad física en el regresor $X$.
\end{opts}

\q{En un modelo semilogarítmico del tipo log-lineal $\ln(Y_i) = \beta_0 + \beta_1 X_i + u_i$ (como la ecuación de ingresos de Mincer), el parámetro $\beta_1$ se interpreta como:}
\begin{opts}
    \item una semielasticidad: un aumento unitario en $X$ cambia $Y$ en aproximadamente $100 \times \beta_1\%$.
    \item el cambio absoluto promedio en $Y$ dividido por el logaritmo neperiano de la media de $X$.
    \item el coeficiente de correlación lineal simple ajustado entre las variables originales $Y$ y $X$.
    \item una elasticidad constante unitaria: ante un cambio de $1\%$ en $X$, $Y$ cambia en $\beta_1$ unidades.
\end{opts}

\q{La presencia de 'Multicolinealidad Imperfecta' o severa en una regresión múltiple tiene como consecuencia econométrica que:}
\begin{opts}
    \item los estimadores MCO pierden insesgadez y consistencia en muestras grandes.
    \item MCO sigue siendo insesgado, pero las varianzas de los estimadores aumentan, reduciendo los estadísticos $t$.
    \item la matriz $X'X$ deja de ser invertible para calcular el vector $\hat{\beta}$.
    \item el coeficiente de determinación $R^2$ cae automáticamente a valores negativos.
\end{opts}

\q{El problema de 'Heterocedasticidad' se define formalmente cuando la varianza condicional de los errores no es constante entre observaciones:}
\begin{opts}
    \item $X'u \neq 0$ en muestras pequeñas con constante incluida.
    \item $E[u_i|X] \neq 0$, variando según la media de cada regresor.
    \item $Cov(u_i, u_j|X) \neq 0$ para todo par de observaciones $i \neq j$.
    \item $Var(u_i|X) = \sigma_i^2$, variando entre observaciones.
\end{opts}

\q{Cuando existe heterocedasticidad en los errores pero se aplica MCO convencional:}
\begin{opts}
    \item la prueba de Durbin-Watson rechaza sistemáticamente la hipótesis nula de perturbaciones esféricas.
    \item el estimador de la varianza residual $s^2$ converge a cero a medida que aumenta el número de variables explicativas.
    \item $\hat{\beta}$ sigue siendo insesgado, pero pierde eficiencia y los errores estándar están sesgados, invalidando los tests $t$ y $F$.
    \item los estimadores $\hat{\beta}$ pierden la insesgadez y la consistencia asintótica incluso bajo exogeneidad estricta.
\end{opts}

\q{La prueba general de White para detectar heterocedasticidad consiste en regresar los residuos al cuadrado $e_i^2$ sobre:}
\begin{opts}
    \item el logaritmo de los errores estandarizados y los valores ajustados al cubo.
    \item todos los regresores originales, sus cuadrados y sus productos cruzados no redundantes ($n R^2 \sim \chi^2_p$).
    \item la variable dependiente rezagada un periodo junto con las variables ficticias.
    \item las diferencias de primer orden de las variables explicativas sin constante.
\end{opts}

\q{La prueba de Breusch-Pagan para heterocedasticidad asume que la varianza del error es una función lineal de un conjunto de variables explicativas $Z_1, \dots, Z_m$. Su hipótesis nula $H_0$ es:}
\begin{opts}
    \item $H_0: \sigma_1^2 = \dots = \sigma_n^2$ (los coeficientes de los regresores $Z$ en la auxiliar de $e_i^2$ son nulos).
    \item $H_0$: los residuos siguen una distribución exponencial negativa con varianza estrictamente decreciente.
    \item $H_0$: el término constante de la regresión original es cero y la media de los residuos supera la varianza.
    \item $H_0$: no existe multicolinealidad severa ni autocorrelación de primer orden entre los regresores $X$.
\end{opts}

\q{El método de Mínimos Cuadrados Ponderados (WLS / MCG) para corregir la heterocedasticidad cuando se conoce la estructura de la varianza $Var(u_i|X) = \sigma^2 h_i$ consiste en:}
\begin{opts}
    \item eliminar el $10\%$ de las observaciones con mayores residuos al cuadrado para evitar apalancamiento extremo.
    \item elevar al cuadrado todas las variables explicativas y estimar la regresión auxiliar sin término constante.
    \item dividir cada observación de todas las variables (incluida la constante) entre $\sqrt{h_i}$ y aplicar MCO.
    \item aplicar primeras diferencias logarítmicas a la variable dependiente para estabilizar la escala residual.
\end{opts}

\q{Los errores estándar robustos a la heterocedasticidad propuestos por White (estimador 'Sandwich' o HC0):}
\begin{opts}
    \item exigen que el tamaño de muestra sea menor a 30 observaciones y corrigen la media de la variable dependiente.
    \item mantienen intactos los $\hat{\beta}$ de MCO pero corrigen la matriz de covarianzas mediante el estimador tipo sándwich.
    \item requieren que las perturbaciones sigan obligatoriamente una distribución $t$ de Student con varianza conocida.
    \item modifican las estimaciones puntuales de los coeficientes $\hat{\beta}$ mediante mínimos cuadrados ponderados en dos etapas.
\end{opts}

\q{El fenómeno de 'Autocorrelación Serial' de primer orden $AR(1)$ en modelos de series de tiempo implica que la perturbación en el periodo $t$ sigue el proceso:}
\begin{opts}
    \item $u_t = \sum_{i=1}^t X_i$, sin término de innovación aleatorio ni memoria.
    \item $u_t = \rho u_{t-1} + \varepsilon_t$, con $|\rho| < 1$ y $\varepsilon_t$ ruido blanco.
    \item $u_t = u_t^2 + \varepsilon_t$, donde la varianza depende linealmente del intercepto.
    \item $u_t = \alpha Y_t + \varepsilon_t$, con $\alpha > 1$ y $\varepsilon_t$ tendencia determinística.
\end{opts}

\q{Cuando existe autocorrelación serial positiva ($\rho > 0$) en los errores de un modelo de series de tiempo estimado por MCO:}
\begin{opts}
    \item el estadístico de Durbin-Watson converge sistemáticamente hacia un valor cercano a 4 en muestras grandes.
    \item los regresores se vuelven perfectamente colineales con sus rezagos, haciendo singular la matriz $(X'X)$.
    \item los estimadores de los coeficientes $\hat{\beta}$ se vuelven sesgados e inconsistentes aun con regresores estrictamente exógenos.
    \item $\hat{\beta}$ sigue insesgado bajo exogeneidad estricta, pero los errores estándar subestiman la varianza, inflando los tests $t$.
\end{opts}

\q{El estadístico $d$ de Durbin-Watson para contrastar autocorrelación serial de primer orden se aproxima algebraicamente a la fórmula $d \approx 2(1 - \hat{\rho})$. Por ende, en ausencia total de autocorrelación ($\hat{\rho} = 0$), el valor del estadístico es:}
\begin{opts}
    \item $d \approx 2$
    \item $d \approx 0$
    \item $d \approx -1$
    \item $d \approx 4$
\end{opts}

\q{Una limitación crucial del estadístico $d$ de Durbin-Watson es que queda estrictamente invalidado si:}
\begin{opts}
    \item los datos provienen de encuestas de corte transversal ordenadas de forma estrictamente aleatoria.
    \item el modelo contiene un término constante junto con más de dos regresores continuos exógenos.
    \item el número de observaciones de la serie supera el umbral de 100 periodos trimestrales.
    \item el modelo incluye la variable dependiente rezagada ($Y_{t-1}$) entre los regresores explicativos.
\end{opts}

\q{La prueba de Breusch-Godfrey (o prueba del Multiplicador de Lagrange - LM) para autocorrelación es superior a la de Durbin-Watson porque:}
\begin{opts}
    \item contrasta autocorrelación de orden superior $AR(p)$ y es válida aun cuando existan rezagos de $Y_t$.
    \item no requiere calcular los residuos de MCO ni estimar una regresión auxiliar sobre los regresores.
    \item se aplica únicamente a modelos con coeficientes estandarizados sin intercepto ni tendencias.
    \item elimina automáticamente los valores atípicos mediante ponderación iterativa de Cochrane-Orcutt.
\end{opts}

\q{Para realizar inferencia estadística válida en series de tiempo en presencia conjunta de heterocedasticidad y autocorrelación de forma desconocida, se utilizan:}
\begin{opts}
    \item los errores estándar HAC consistentes de Newey y West (1987).
    \item los contrastes de raíz unitaria aumentados de Dickey-Fuller (ADF).
    \item el filtro cíclico de Hodrick-Prescott aplicado sobre los regresores.
    \item los errores estándar clásicos de MCO ajustados por grados de libertad.
\end{opts}

\q{El problema de 'Endogeneidad' se produce cuando uno o más regresores están correlacionados con la perturbación estocástica ($Cov(X, u) \neq 0$). Sus tres causas más comunes en economía son:}
\begin{opts}
    \item transformación logarítmica de la variable dependiente $Y$.
    \item homocedasticidad condicional, normalidad y estacionariedad.
    \item sesgo de variable omitida, error de medición en $X$ y simultaneidad.
    \item ausencia de constante, colinealidad exacta y exceso de muestra.
\end{opts}

\q{La consecuencia econométrica más grave de la endogeneidad en el estimador de MCO es que:}
\begin{opts}
    \item el estimador $\hat{\beta}$ es sesgado e inconsistente ($plim \, \hat{\beta} \neq \beta$).
    \item el coeficiente de correlación de Pearson entre los regresores se vuelve exactamente igual a 1.
    \item el estimador MCO conserva su consistencia asintótica pero pierde eficiencia en muestras grandes.
    \item los errores estándar calculados convergen a cero provocando un sobreajuste de los residuos.
\end{opts}

\q{Para resolver la endogeneidad mediante el método de Variables Instrumentales (VI / IV), una variable instrumental válida $Z$ debe satisfacer dos condiciones indispensables:}
\begin{opts}
    \item estar altamente correlacionada con la perturbación $u$ y no correlacionada con $X$.
    \item relevancia del instrumento ($Cov(Z, X) \neq 0$) y exogeneidad ($Cov(Z, u) = 0$).
    \item seguir una distribución uniforme continua y ser ortogonal a la variable explicativa $X$.
    \item ser una transformación cuadrática exacta de la variable dependiente observada en el modelo.
\end{opts}

\q{El estimador de Mínimos Cuadrados en Dos Etapas (MC2E / 2SLS) resuelve la endogeneidad siguiendo el procedimiento:}
\begin{opts}
    \item Etapa 1: Diferenciar las series temporales; Etapa 2: Aplicar mínimos cuadrados ponderados sobre los residuos.
    \item Etapa 1: Regresar $Y$ sobre los instrumentos $Z$; Etapa 2: Regresar la variable endógena $X$ sobre los residuos.
    \item Etapa 1: Regresar $X$ sobre los instrumentos $Z$ para obtener $\hat{X}$; Etapa 2: Regresar $Y$ sobre $\hat{X}$ por MCO.
    \item Etapa 1: Calcular la matriz de correlaciones parciales; Etapa 2: Eliminar las variables con menor estadístico $t$.
\end{opts}


\section*{Bloque B. Aplicaciones cuantitativas y contrastes econométricos}

\noindent\textit{Preguntas 31 a 65. Evalúan la capacidad de cálculo analítico de estimadores, varianzas residuales, pruebas de significancia, diagnósticos de heterocedasticidad (White, Breusch-Pagan), autocorrelación (Durbin-Watson, Prais-Winsten), sesgo de variable omitida, variables instrumentales (2SLS) y pruebas de especificación.}

\q{En una regresión lineal simple $Y_i = \beta_0 + \beta_1 X_i + u_i$ con $n = 50$ observaciones, se tiene: $\sum (X_i - \bar{X})^2 = 400$, $\sum (Y_i - \bar{Y})^2 = 900$ y $\sum (X_i - \bar{X})(Y_i - \bar{Y}) = 600$. ¿Cuál es el valor estimado de la pendiente MCO $\hat{\beta}_1$?}
\begin{opts}
    \item $\hat{\beta}_1 = 1.50$
    \item $\hat{\beta}_1 = 1.00$
    \item $\hat{\beta}_1 = 0.67$
    \item $\hat{\beta}_1 = 2.25$
\end{opts}

\q{Para los datos de la pregunta anterior, determine el coeficiente de correlación lineal de Pearson $r$ y el coeficiente de determinación $R^2$ del modelo:}
\begin{opts}
    \item $r = 0.75$ y $R^2 = 0.56$
    \item $r = 0.50$ y $R^2 = 0.25$
    \item $r = 0.80$ y $R^2 = 0.64$
    \item $r = 1.00$ y $R^2 = 1.00$
\end{opts}

\q{En un modelo de regresión con $n = 45$ observaciones y $k = 5$ parámetros estimados, la Suma de Cuadrados Residuales es $SCR = 800$. ¿Cuál es la varianza residual estimada ($s^2$) y el error estándar de la regresión ($s$)?}
\begin{opts}
    \item $s^2 = 17.78$ y $s = 4.22$
    \item $s^2 = 16$ y $s = 4.00$
    \item $s^2 = 20$ y $s = 4.47$
    \item $s^2 = 25$ y $s = 5.00$
\end{opts}

\q{En una regresión de demanda de agua, el coeficiente estimado de la variable precio es $\hat{\beta}_1 = -0.45$ con un error estándar de $EE(\hat{\beta}_1) = 0.15$. El tamaño muestral es $n = 120$ y el modelo tiene $k = 4$ parámetros. Para contrastar $H_0: \beta_1 = 0$ frente a $H_1: \beta_1 \neq 0$ con un nivel de significancia del $5\%$ (valor crítico bilateral $t_{0.025, 116} \approx 1.98$):}
\begin{opts}
    \item El estadístico es $t = -3.00$, por lo que se rechaza $H_0$ al $5\%$.
    \item El estadístico es $t = -1.50$, por lo que no se rechaza $H_0$ al $5\%$.
    \item El estadístico es $t = -0.067$, resultado no significativo al $5\%$.
    \item El estadístico es $t = -6.75$, resultado no significativo al $5\%$.
\end{opts}

\q{Un modelo de regresión múltiple con $k = 4$ parámetros (3 variables explicativas más constante) y $n = 64$ observaciones arroja un coeficiente de determinación de $R^2 = 0.60$. ¿Cuál es el valor del estadístico $F$ de significancia global de la regresión?}
\begin{opts}
    \item $F = 30.0$
    \item $F = 15.0$
    \item $F = 45.0$
    \item $F = 20.0$
\end{opts}

\q{Se contrastan 2 restricciones lineales simultáneas ($q = 2$) en un modelo no restringido con $k = 6$ parámetros y $n = 106$ observaciones. La suma de cuadrados residuales no restringida es $SCR_{NR} = 500$ y la restringida es $SCR_R = 560$. El valor del estadístico $F$ es:}
\begin{opts}
    \item $F = 8.0$
    \item $F = 15.0$
    \item $F = 6.0$
    \item $F = 12.0$
\end{opts}

\q{La estimación de una función de demanda residencial de agua arroja el modelo: $\ln(Q) = 2.4 - 0.35 \ln(P) + 0.40 \ln(Y)$, donde $Q$ es el consumo mensual en $m^3$, $P$ es la tarifa por $m^3$ y $Y$ es el ingreso familiar. Si la tarifa de agua sube en un $10\%$, manteniendo el ingreso constante:}
\begin{opts}
    \item el consumo de agua se reduce en $3.5$ metros cúbicos exactos.
    \item el consumo de agua se reduce en aproximadamente un $3.5\%$.
    \item el consumo de agua aumenta en $4.0\%$ por efecto sustitución.
    \item la elasticidad precio de la demanda cambia a $-0.035$.
\end{opts}

\q{En una ecuación semilogarítmica de costos operativos de EPS: $\ln(OPEX) = 8.2 + 0.04 \, Conexiones$, donde $Conexiones$ está expresada en miles de conexiones activas. Si una EPS añade 5 mil conexiones adicionales a su red:}
\begin{opts}
    \item sus costos operativos aumentan en aproximadamente $20\%$.
    \item sus costos operativos aumentan en un $4\%$ acumulado.
    \item la elasticidad costo de las conexiones es igual a $0.04$.
    \item sus costos operativos aumentan en 200 soles mensuales.
\end{opts}

\q{En un modelo lineal-log $Y = 50 + 200 \ln(X)$, si la variable $X$ se incrementa en un $1\%$:}
\begin{opts}
    \item la variable dependiente $Y$ aumenta en $2$ unidades.
    \item la variable dependiente $Y$ aumenta en un $1\%$ exacto.
    \item la variable dependiente $Y$ aumenta en $200$ unidades.
    \item la variable dependiente $Y$ se duplica proporcionalmente.
\end{opts}

\q{Al analizar la multicolinealidad entre dos variables explicativas en un modelo de costos de saneamiento, la regresión auxiliar de una de ellas sobre las demás arroja un coeficiente de determinación de $R_j^2 = 0.80$. ¿Cuál es el Factor de Inflación de la Varianza ($VIF$) correspondiente?}
\begin{opts}
    \item $VIF = 1.25$
    \item $VIF = 5.0$
    \item $VIF = 10.0$
    \item $VIF = 4.0$
\end{opts}

\q{En un modelo con dos regresores y constante, si la matriz $(X'X)^{-1}$ tiene en su diagonal principal los elementos $[0.5, 0.04, 0.09]$ y la varianza residual estimada es $s^2 = 25$, ¿cuál es el error estándar del segundo estimador de pendiente $EE(\hat{\beta}_2)$?}
\begin{opts}
    \item $EE(\hat{\beta}_2) = 1.50$
    \item $EE(\hat{\beta}_2) = 5.00$
    \item $EE(\hat{\beta}_2) = 2.25$
    \item $EE(\hat{\beta}_2) = 0.75$
\end{opts}

\q{Para contrastar heterocedasticidad mediante la prueba de White en un modelo con 2 regresores sin interacciones (la regresión auxiliar incluye constante, $X_1, X_2, X_1^2, X_2^2, X_1 X_2$, total 5 regresores excluyendo constante), se tiene $n = 100$ y el $R^2$ auxiliar es $R_{aux}^2 = 0.15$. El estadístico $LM$ de White y su distribución son:}
\begin{opts}
    \item $LM = 1.50 \sim t_{95}$
    \item $LM = 15.0 \sim \chi^2_5$
    \item $LM = 30.0 \sim F_{5, 95}$
    \item $LM = 7.5 \sim \chi^2_2$
\end{opts}

\q{En una prueba de Breusch-Pagan con $m = 3$ variables explicativas en la ecuación de la varianza y $n = 80$ observaciones, el $R^2$ de la regresión auxiliar de residuos al cuadrado es $R_{BP}^2 = 0.10$. El estadístico $LM = n R_{BP}^2$ es:}
\begin{opts}
    \item $LM = 2.4$
    \item $LM = 0.8$
    \item $LM = 26.67$
    \item $LM = 8.0$
\end{opts}

\q{Si se conoce que la varianza del error en un modelo de corte transversal es proporcional al cuadrado de una variable $X_i$ ($Var(u_i|X) = \sigma^2 X_i^2$), la transformación adecuada de Mínimos Cuadrados Ponderados (WLS) sobre la ecuación $Y_i = \beta_0 + \beta_1 X_i + u_i$ exige estimar:}
\begin{opts}
    \item $\frac{Y_i}{X_i^2} = \frac{\beta_0}{X_i^2} + \frac{\beta_1}{X_i} + \frac{u_i}{X_i^2}$
    \item $\ln(Y_i) = \beta_0 + \beta_1 \ln(X_i) + u_i$
    \item $\frac{Y_i}{X_i} = \beta_0 \left(\frac{1}{X_i}\right) + \beta_1 + \frac{u_i}{X_i}$
    \item $Y_i X_i = \beta_0 X_i + \beta_1 X_i^2 + u_i X_i$
\end{opts}

\q{En una regresión de series de tiempo trimestrales de consumo de agua con $n = 100$ observaciones, el coeficiente de autocorrelación muestral de primer orden de los residuos es $\hat{\rho} = 0.60$. ¿Cuál es el valor aproximado del estadístico de Durbin-Watson $d$?}
\begin{opts}
    \item $d \approx 2.00$
    \item $d \approx 1.20$
    \item $d \approx 3.20$
    \item $d \approx 0.80$
\end{opts}

\q{Si los límites críticos de Durbin-Watson para $n = 100$ y $k' = 2$ regresores al nivel del $5\%$ son $d_L = 1.63$ y $d_U = 1.72$, y el valor calculado es $d = 0.80$, la conclusión estadística es:}
\begin{opts}
    \item no existe evidencia de autocorrelación serial en el error.
    \item la prueba resulta no concluyente por zona de indecisión.
    \item existe evidencia de autocorrelación serial negativa.
    \item existe evidencia de autocorrelación serial positiva.
\end{opts}

\q{Si en un modelo de series de tiempo el estadístico de Durbin-Watson resulta ser $d = 1.30$, ¿cuál es el valor estimado del parámetro autorregresivo $\hat{\rho}$ para aplicar la transformación de Cochrane-Orcutt?}
\begin{opts}
    \item $\hat{\rho} \approx 0.35$
    \item $\hat{\rho} \approx -0.35$
    \item $\hat{\rho} \approx 0.15$
    \item $\hat{\rho} \approx 0.70$
\end{opts}

\q{Bajo la transformación cuasi-diferenciada de Prais-Winsten con parámetro $\hat{\rho} = 0.35$, la variable dependiente transformada para el primer periodo ($t=1$) se define como:}
\begin{opts}
    \item $Y_1^* = (1 - 0.35^2) \, Y_1$
    \item $Y_1^* = \sqrt{1 - (0.35)^2} \, Y_1$
    \item $Y_1^* = (1 - 0.35^2)^{-1/2} \, Y_1$
    \item $Y_1^* = Y_1 - 0.35 Y_0$
\end{opts}

\q{El verdadero modelo poblacional es $Y = 10 + 4 X_1 + 2 X_2 + u$, pero el investigador omite por error la variable $X_2$ y estima por MCO la regresión simple $Y = \tilde{\beta}_0 + \tilde{\beta}_1 X_1 + v$. Si la regresión auxiliar de la variable omitida es $X_2 = 1.5 + 0.5 X_1$, ¿a qué valor convergerá el estimador de la pendiente $\tilde{\beta}_1$?}
\begin{opts}
    \item $\tilde{\beta}_1$ converge a 4.0
    \item $\tilde{\beta}_1$ converge a 2.0
    \item $\tilde{\beta}_1$ converge a 5.0
    \item $\tilde{\beta}_1$ converge a 6.0
\end{opts}

\q{En una regresión simple con regresor endógeno $X$, se utiliza un instrumento exógeno $Z$. Se dispone de: $\widehat{Cov}(Z, Y) = 150$ y $\widehat{Cov}(Z, X) = 75$. ¿Cuál es el estimador de Variables Instrumentales $\hat{\beta}_{IV}$?}
\begin{opts}
    \item $\hat{\beta}_{IV} = 1.50$
    \item $\hat{\beta}_{IV} = 2.00$
    \item $\hat{\beta}_{IV} = 11250$
    \item $\hat{\beta}_{IV} = 0.50$
\end{opts}

\q{En la estimación por Mínimos Cuadrados en Dos Etapas (MC2E / 2SLS) con un regresor endógeno, la regresión de primera etapa arrojó un estadístico $F$ sobre los instrumentos excluidos de $F = 4.2$. Siguiendo la regla de oro de Stock, Wright y Yogo (2002) para detectar 'Instrumentos Débiles':}
\begin{opts}
    \item los instrumentos son débiles ($F < 10$), sesgando el estimador 2SLS hacia MCO.
    \item los instrumentos superan el umbral y garantizan inferencia no sesgada.
    \item la prueba de Durbin-Watson rechaza la existencia de endogeneidad.
    \item el modelo es exactamente identificado y debe estimarse por MCO.
\end{opts}

\q{Para verificar si una variable sospechosa $X_1$ es realmente endógena mediante la prueba de especificación de Durbin-Wu-Hausman en forma de regresión:}
\begin{opts}
    \item se añaden los residuos de la primera etapa $\hat{v}$ en la ecuación MCO y se evalúa su prueba $t$.
    \item se calcula la correlación simple entre $X_1$ y la constante tras aplicar primeras diferencias al panel.
    \item se reemplaza $X_1$ por su inversa y se contrasta si el intercepto de la nueva ecuación es nulo al $5\%$.
    \item se estima el modelo por MCO y por WLS y se exige que la diferencia entre ambos $R^2$ supere el $90\%$.
\end{opts}

\q{En un modelo con 1 variable endógena y 3 instrumentos excluidos (modelo sobreidentificado con $m = 3$ y $k_1 = 1$, grados de sobreidentificación $m - k_1 = 2$), la prueba de restricciones de sobreidentificación de Sargan arroja un estadístico de $Sargan = n R_{aux}^2 = 1.20$. Si el valor crítico al $5\%$ de $\chi^2_2$ es $5.99$, la conclusión es:}
\begin{opts}
    \item el modelo está subidentificado y requiere añadir más variables endógenas.
    \item la prueba indica que los instrumentos son débiles en la primera etapa.
    \item no se rechaza $H_0$: los instrumentos cumplen la restricción de exclusión.
    \item se rechaza $H_0$: los instrumentos son inválidos y correlacionados con $u$.
\end{opts}

\q{Al modelar diferencias de costos operativos entre regiones con 4 categorías (Costa, Sierra, Selva y Lima Metropolitana), si un economista incluye una variable ficticia (dummy) para cada una de las 4 regiones y además incluye el término constante en la regresión MCO:}
\begin{opts}
    \item se incurre en la Trampa de Variables Ficticias, generando multicolinealidad perfecta.
    \item el modelo gana máxima precisión al capturar efectos fijos completos de las cuatro zonas.
    \item los coeficientes miden directamente la desviación estándar de cada categoría regional.
    \item el $R^2$ ajustado aumenta artificialmente al cuadruplicar los grados de libertad.
\end{opts}

\q{En un modelo salarial $\ln(W) = \beta_0 + \beta_1 Educ + \beta_2 Mujer + \beta_3 (Educ \times Mujer) + u$, donde $Mujer$ es una dummy (1 si es mujer, 0 si es hombre), el retorno porcentual a un año adicional de educación para una mujer es:}
\begin{opts}
    \item $\beta_1$
    \item $\beta_1 \times \beta_3$
    \item $\beta_2 + \beta_3$
    \item $\beta_1 + \beta_3$
\end{opts}

\q{En una regresión lineal con variables estandarizadas (donde cada variable se transforma restando su media y dividiendo entre su desviación estándar $Z_X = \frac{X - \bar{X}}{S_X}$):}
\begin{opts}
    \item la varianza residual estimada es exactamente igual a $n - k$.
    \item el $R^2$ se convierte en un número estrictamente negativo.
    \item el intercepto estimado $\hat{\beta}_0$ es idénticamente cero.
    \item los errores estándar de todas las pendientes se duplican.
\end{opts}

\q{La matriz de proyección o 'Matriz Sombrero' (Hat Matrix) definida como $H = X(X'X)^{-1}X'$ cumple las propiedades algebraicas de ser:}
\begin{opts}
    \item asimétrica y no invertible bajo cualquier rango de $X$.
    \item simétrica ($H' = H$) e idempotente ($H^2 = H$).
    \item antisimétrica, definida negativa y ortogonal a $X$.
    \item estrictamente diagonal con ceros en la diagonal principal.
\end{opts}

\q{Una propiedad algebraica exacta de los residuos de MCO ($e = Y - X\hat{\beta}$) en cualquier modelo que incluya término constante es que:}
\begin{opts}
    \item la media de los residuos coincide con la varianza incondicional de $Y$.
    \item el producto de los residuos individuales es siempre igual a la unidad ($1$).
    \item la suma de los residuos es cero ($\sum e_i = 0$) y se cumple $X'e = 0$.
    \item los residuos son estocásticamente independientes entre sí para todo tamaño $n$.
\end{opts}

\q{La descomposición clásica de la varianza en MCO establece que la Suma Total de Cuadrados ($STC$) se descompone en Suma de Cuadrados Explicada ($SEC$) más Suma de Cuadrados Residuales ($SCR$): $STC = SEC + SCR$. Esta igualdad se cumple estrictamente siempre que:}
\begin{opts}
    \item el modelo lineal incluya un término constante (intercepto).
    \item todos los regresores sigan una distribución normal bivariada.
    \item el tamaño muestral exceda las 1000 observaciones continuas.
    \item no exista correlación lineal alguna entre $Y$ y $X$.
\end{opts}

\q{Para un coeficiente estimado $\hat{\beta}_1 = 2.50$ con $EE(\hat{\beta}_1) = 0.25$ en una muestra con $n - k = 60$ grados de libertad (valor crítico $t_{0.025} = 2.00$), el intervalo de confianza al $95\%$ es:}
\begin{opts}
    \item $[2.45, \; 2.55]$
    \item $[2.25, \; 2.75]$
    \item $[2.00, \; 3.00]$
    \item $[1.50, \; 3.50]$
\end{opts}

\q{Al realizar una predicción individual puntual de $Y_0$ para un vector de valores de regresores $x_0$, la varianza del error de predicción individual $Var(Y_0 - \hat{Y}_0)$ es:}
\begin{opts}
    \item $\sigma^2 [x_0' (X'X)^{-1} x_0]$
    \item $\sigma^2 [n + x_0' (X'X) x_0]$
    \item $\sigma^2 [1 + x_0' (X'X)^{-1} x_0]$
    \item $\sigma^2 [1 - x_0' (X'X)^{-1} x_0]$
\end{opts}

\q{En la selección entre especificaciones de modelos rivales, el Criterio de Información de Schwarz (BIC / SIC) se diferencia del Criterio de Akaike (AIC) en que:}
\begin{opts}
    \item el AIC solo se utiliza en modelos de variables instrumentales y panel dinámico.
    \item el BIC penaliza más parámetros adicionales ($\ln(n)k$ vs $2k$ en AIC).
    \item el BIC exige que el $R^2$ ajustado supere $0.80$ para seleccionar modelos.
    \item el BIC no incorpora ningún término de penalización por grados de libertad.
\end{opts}

\q{La prueba de normalidad de Jarque-Bera sobre los residuos de MCO evalúa si el coeficiente de asimetría muestral $S$ y la curtosis muestral $K$ coinciden con los de una distribución normal teórica ($S=0, K=3$). Su estadístico es:}
\begin{opts}
    \item $JB = \frac{n}{24} \left( S^2 + \frac{K^2}{4} \right) \sim \chi^2_4$
    \item $JB = \frac{n}{6} \left( S^2 + (K - 3)^2 \right) \sim \chi^2_2$
    \item $JB = \frac{n}{12} \left( (S - 1)^2 + \frac{(K - 3)^2}{2} \right) \sim F_{2, n}$
    \item $JB = \frac{n}{6} \left( S^2 + \frac{(K - 3)^2}{4} \right) \sim \chi^2_2$
\end{opts}

\q{La prueba RESET de Ramsey (Regression Equation Specification Error Test) se utiliza en econometría para:}
\begin{opts}
    \item determinar el número óptimo de rezagos autorregresivos $p$ y de media móvil $q$ en series ARIMA.
    \item detectar errores de forma funcional añadiendo potencias $\hat{Y}^2, \hat{Y}^3$ con test $F$.
    \item contrastar si existe heterocedasticidad condicional autorregresiva en los residuos de la muestra.
    \item calcular la elasticidad precio cruzada en sistemas de ecuaciones aparentemente no relacionadas.
\end{opts}

\q{Al estimar un modelo de MCO con restricciones lineales $R\beta = r$ (por ejemplo, rendimientos constantes a escala en una función Cobb-Douglas $\beta_1 + \beta_2 = 1$), el estimador MCO restringido $\hat{\beta}_R$ frente al no restringido $\hat{\beta}_{NR}$:}
\begin{opts}
    \item impide calcular matriz de covarianzas o errores estándar robustos de White.
    \item siempre presenta $SCR_R \le SCR_{NR}$, logrando un mayor ajuste de $R^2$.
    \item se vuelve inconsistente asintóticamente aunque las restricciones sean ciertas.
    \item siempre presenta $SCR_R \ge SCR_{NR}$, al optimizar en un espacio acotado.
\end{opts}


\section*{Bloque C. Casos con implicancias regulatorias y sector saneamiento}

\noindent\textit{Preguntas 66 a 90. Evalúan aplicaciones econométricas avanzadas al sector de agua y saneamiento: estimación de elasticidades de demanda, tarifas por bloques crecientes, fronteras estocásticas de costos (SFA) en EPS, diferencias en diferencias (Diff-in-Diff), regresión discontinua en SISFOH, WACC y modelos de datos de panel.}

\q{En la estimación econométrica de la demanda de agua potable residencial para fijar tarifas en SUNASS, la especificación comúnmente adoptada es de tipo log-log: $\ln(Q_i) = \alpha + \beta_P \ln(P_i) + \beta_Y \ln(Y_i) + \gamma Z_i + u_i$. Los parámetros $\beta_P$ y $\beta_Y$ representan directamente:}
\begin{opts}
    \item el consumo básico de subsistencia y el cargo fijo mensual.
    \item la elasticidad precio directa y la elasticidad ingreso de la demanda.
    \item la semielasticidad tarifaria media y el costo marginal.
    \item el margen operativo regulado y el factor de productividad.
\end{opts}

\q{En los sistemas de tarifas de agua por bloques crecientes (Increasing Block Tariffs - IBT) utilizados por las EPS peruanas, la estimación por MCO de la demanda de agua enfrenta un grave problema de endogeneidad porque:}
\begin{opts}
    \item el agua potable es un bien homogéneo cuya utilidad marginal decreciente impide estimar regresiones.
    \item el precio marginal y la cantidad consumida se determinan simultáneamente por el salto de bloque.
    \item la Superintendencia no publica los pliegos tarifarios vigentes de cada bloque en el diario El Peruano.
    \item el consumo de agua residencial es estacionalmente fijo e invariante ante cambios tarifarios mensuales.
\end{opts}

\q{Para corregir la endogeneidad de la tarifa marginal en modelos de demanda de agua con bloques crecientes, la literatura econométrica y los estudios regulatorios emplean como variable instrumental:}
\begin{opts}
    \item la tarifa del primer bloque o estructura tarifaria exógena de EPS comparables.
    \item la tasa de interés de referencia del BCRP y el encaje bancario comercial promedio.
    \item el índice de precios al consumidor de alimentos y bebidas del área metropolitana.
    \item el producto bruto interno per cápita nacional ajustado por paridad cambiaria.
\end{opts}

\q{SUNASS utiliza modelos de Análisis de Frontera Estocástica (Stochastic Frontier Analysis - SFA) para estimar los costos eficientes de las EPS. La ventaja econométrica clave de SFA frente a una regresión MCO estándar es que:}
\begin{opts}
    \item asume que todas las empresas prestadoras operan exactamente sobre la frontera con retorno del $100\%$.
    \item descompone el error en ineficiencia técnica no negativa (asimétrica) y ruido aleatorio simétrico.
    \item elimina la necesidad de especificar forma funcional de costos o calcular pruebas de significancia $t$.
    \item no requiere recopilar información contable ni estados financieros auditados de las empresas EPS.
\end{opts}

\q{Al estimar una función de costos translogarítmica multirregulada para las 50 EPS del Perú, se evalúan simultáneamente economías de escala y economías de densidad de clientes. La economía de densidad se define econométricamente como:}
\begin{opts}
    \item la tasa de depreciación de los micromedidores instalados por kilómetro.
    \item la variación porcentual del costo ante mayor volumen manteniendo fijos red y clientes.
    \item la reducción del costo medio cuando se fusionan dos EPS distantes.
    \item la relación técnica entre el gasto laboral y el costo de capital.
\end{opts}

\q{En una regresión de corte transversal de costos de operación (OPEX) que incluye a Sedapal junto a empresas provinciales pequeñas (como EPS Moyobamba o EPS Ilo), la presencia de heterocedasticidad es sumamente probable debido a que:}
\begin{opts}
    \item las empresas comparadas utilizan tecnologías distintas de micromedición ultrasónica.
    \item los residuos de la regresión siguen necesariamente una distribución uniforme discreta.
    \item todas las EPS aplican exactamente la misma estructura tarifaria para uso doméstico.
    \item la enorme diferencia de escala y dispersión entre Sedapal y las pequeñas EPS locales.
\end{opts}

\q{Para corregir la distorsión inferencial causada por la heterocedasticidad en la estimación de costos de EPS con muestras pequeñas a medianas ($n \approx 50$), la literatura econométrica recomienda utilizar la corrección de White modificada:}
\begin{opts}
    \item estimador sándwich asintótico HC0 sin ajuste de muestra.
    \item regresión de cuantiles no condicionada por mínimos desvíos.
    \item errores estándar convencionales de MCO homocedásticos.
    \item HC3 (Jackknife), que penaliza observaciones con alto leverage.
\end{opts}

\q{Al modelar series de tiempo mensuales de caudales de captación de agua en ríos de la vertiente del Pacífico para planificar el abastecimiento de EPS, la presencia de estacionalidad y persistencia intertemporal se modela mediante:}
\begin{opts}
    \item modelos SARIMA, que capturan autocorrelación serial y ciclos estacionales.
    \item análisis de componentes principales no lineales sobre series anuales.
    \item regresión lineal simple con intercepto nulo y tendencia determinística.
    \item modelos de selección muestral en dos etapas de corte transversal.
\end{opts}

\q{En modelos de series de tiempo de demanda de agua urbana donde se evalúa el impacto de episodios de El Niño, la aplicación de errores estándar HAC de Newey-West con $L$ rezagos permite:}
\begin{opts}
    \item garantizar que el estadístico de Durbin-Watson sea exactamente igual a 2.
    \item obtener inferencia robusta frente a correlación serial de choques y heterocedasticidad estacional.
    \item sustituir las series históricas pluviométricas de Senamhi mediante interpolación.
    \item eliminar por completo la tendencia poblacional y las variables ficticias.
\end{opts}

\q{En un panel de seguimiento de consumo de agua con micromedidores domiciliarios a lo largo de 5 años, el hecho de que los hogares con medidores defectuosos o manipulados sean retirados sistemáticamente de la base de datos genera un problema de:}
\begin{opts}
    \item violación del Teorema de Gauss-Markov por exceso de grados de libertad en la muestra.
    \item sesgo de atrición no aleatorio (Attrition Bias) correlacionado con el volumen consumido.
    \item multicolinealidad perfecta entre los efectos fijos temporales y la estructura tarifaria.
    \item discriminación de precios de segundo grado en la estimación de elasticidades cruzadas.
\end{opts}

\q{Para evaluar el impacto de un programa piloto de instalación masiva de micromedidores sobre el consumo de agua ejecutado por una EPS, la metodología cuasiexperimental estándar es:}
\begin{opts}
    \item aplicar la prueba de normalidad asintótica de Jarque-Bera sobre la distribución del consumo.
    \item comparar únicamente el promedio de consumo antes y después de la obra sin grupo de comparación.
    \item el estimador de Diferencias en Diferencias comparando grupos con y sin nuevos medidores.
    \item estimar una regresión lineal simple de corte transversal sin incluir término constante.
\end{opts}

\q{El supuesto de identificación indispensable para la validez causal del estimador de Diferencias en Diferencias en la evaluación del programa de micromedición es:}
\begin{opts}
    \item ausencia absoluta de correlación contemporánea entre las características de los hogares.
    \item que el coeficiente de determinación $R^2$ de la regresión auxiliar sea exactamente igual a 1.
    \item que la tarifa marginal media del agua potable iguale el costo medio de provisión del servicio.
    \item tendencias paralelas: sin intervención, tratados y controles habrían seguido igual trayectoria.
\end{opts}

\q{Para medir el impacto de la tarifa subsidiada del SISFOH sobre el bienestar de los hogares pobres, cuando el subsidio se asigna estrictamente si el puntaje de pobreza del hogar es menor a un umbral $c_0$, el método econométrico óptimo es:}
\begin{opts}
    \item Regresión Discontinua (RDD) nítida o difusa alrededor del umbral $c_0$.
    \item análisis de componentes principales de Hotelling sobre variables continuas.
    \item modelo de vectores autorregresivos (VAR) con restricciones de corto plazo.
    \item prueba de autocorrelación de Durbin-Watson para datos de corte transversal.
\end{opts}

\q{Al modelar encuestas de SUNASS sobre la decisión discreta de un hogar de pagar o no una tarifa adicional por agua con 24 horas continuas de servicio ($Y_i \in \{0, 1\}$), la estimación por el Modelo Lineal de Probabilidad (MLP) mediante MCO presenta la debilidad de que:}
\begin{opts}
    \item impide por completo incluir variables ficticias regionales o calcular correlaciones lineales.
    \item predice probabilidades fuera de $[0, 1]$ y sufre heterocedasticidad, prefiriéndose Logit/Probit.
    \item exige un tamaño de muestra infinito para poder estimar los parámetros mediante matrices.
    \item fuerza a que todos los coeficientes estimados converjan a cero cuando hay constante.
\end{opts}

\q{En un modelo de costos de inversión de las EPS que incluye simultáneamente como regresores: longitud de redes ($Km$), número de conexiones de agua ($Conex$) y volumen total producido ($m^3$), es muy probable encontrar Factores de Inflación de Varianza ($VIF > 10$) debido a:}
\begin{opts}
    \item la presencia de subsidios cruzados entre categorías comerciales e industriales.
    \item la correlación serial negativa entre las observaciones de distintas regiones naturales.
    \item la falta de normalidad asintótica de los residuos en la prueba de heterocedasticidad.
    \item la fuerte multicolinealidad estructural de variables que miden escala operativa.
\end{opts}

\q{Al modelar los costos de mantenimiento de redes de una EPS, si se incluye el porcentaje de Agua No Facturada ($ANF$) como regresor explicativo, existe endogeneidad por causalidad inversa porque:}
\begin{opts}
    \item mayor ANF eleva costos de reparación, pero menor mantenimiento preventivo dispara fugas y ANF.
    \item el porcentaje de agua no facturada es un parámetro constante fijado por ley.
    \item no existen registros históricos estandarizados de agua no facturada en las EPS.
    \item el indicador ANF es una variable cualitativa imposible de medir con precisión.
\end{opts}

\q{En encuestas de hogares donde se evalúa el gasto en agua potable, muchas familias periurbanas carecen de conexión a la red pública y no reportan tarifa de EPS. Modelar el gasto en agua solo para los hogares conectados genera un problema econométrico de:}
\begin{opts}
    \item violación del supuesto de homocedasticidad por exceso de grados de libertad.
    \item autocorrelación de orden 12 entre los residuos de consumo mensual domiciliario.
    \item multicolinealidad perfecta entre la variable dependiente y la constante.
    \item sesgo de selección muestral, corregible mediante el procedimiento de Heckman.
\end{opts}

\q{Al estimar la productividad de las 50 EPS mediante un panel de datos durante el periodo regulatorio 2020-2025 ($Y_{it} = \alpha_i + X_{it}\beta + u_{it}$), el modelo de Efectos Fijos (Fixed Effects) es preferible al de MCO agrupado (Pooled) porque:}
\begin{opts}
    \item no requiere descontar grados de libertad al estimar los parámetros de pendiente del modelo.
    \item elimina automáticamente toda heterocedasticidad y autocorrelación serial en los errores.
    \item controla la heterogeneidad individual invariante en el tiempo aun correlacionada con $X_{it}$.
    \item asume que todas las empresas prestadoras comparten exactamente el mismo intercepto individual.
\end{opts}

\q{La prueba de Hausman en modelos de datos de panel compara el estimador de Efectos Fijos ($\hat{\beta}_{FE}$) con el de Efectos Aleatorios ($\hat{\beta}_{RE}$). Si la hipótesis nula $H_0: Cov(\alpha_i, X_{it}) = 0$ es rechazada estadísticamente ($p < 0.05$):}
\begin{opts}
    \item Efectos Aleatorios es inconsistente por correlación con $\alpha_i$, debiendo usarse Efectos Fijos.
    \item ambos estimadores resultan algebraicamente idénticos al estimador MCO agrupado (Pooled).
    \item el estimador de Efectos Aleatorios sigue siendo consistente y supera en eficiencia a Efectos Fijos.
    \item el modelo de panel queda subidentificado y no puede estimarse por método alguno.
\end{opts}

\q{En el cálculo del Costo Promedio Ponderado de Capital (WACC regulatorio) que fija SUNASS para remunerar las inversiones de las EPS, el costo del capital accionario ($K_e$) se estima habitualmente mediante el modelo Capital Asset Pricing Model (CAPM): $K_e = R_f + \beta_L (R_m - R_f) + RP$. En esta ecuación, el coeficiente $\beta_L$ se estima econométricamente mediante:}
\begin{opts}
    \item el logaritmo natural de la razón entre volumen facturado y volumen producido por Sedapal.
    \item el cociente contable promedio entre el activo fijo neto y el pasivo corriente de la EPS.
    \item la prueba de Durbin-Watson aplicada sobre la serie mensual de tasa de interés interbancaria.
    \item regresión de retornos de acciones reguladas comparables sobre el retorno del mercado.
\end{opts}

\q{Dado que las EPS peruanas son empresas estatales no cotizadas en bolsa, SUNASS utiliza Betas de empresas comparables internacionales (proxy) desapalancadas y reapalancadas mediante la fórmula de Hamada: $\beta_L = \beta_U [1 + (1 - T)(D/E)]$. Si la Beta desapalancada promedio es $\beta_U = 0.50$, la tasa del impuesto a la renta es $T = 30\%$ y la estructura óptima de endeudamiento es $D/E = 0.60$, ¿cuál es la Beta apalancada regulada $\beta_L$?}
\begin{opts}
    \item $\beta_L = 1.00$
    \item $\beta_L = 0.85$
    \item $\beta_L = 0.50$
    \item $\beta_L = 0.71$
\end{opts}

\q{Para modelar econométricamente el número de roturas de tuberías o incidentes de atoros de alcantarillado por sector hidrográfico ($Y_i = 0, 1, 2, \dots$, datos de conteo), el modelo más apropiado que supera la limitación de la regresión lineal es:}
\begin{opts}
    \item el modelo de paseo aleatorio con deriva y raíz unitaria estacional.
    \item la prueba de cointegración en dos etapas de Engle y Granger.
    \item regresión de Poisson o Binomial Negativo (si hay sobredispersión).
    \item el modelo de Mínimos Cuadrados Ordinarios simple sin intercepto.
\end{opts}

\q{En la detección forense de colusión en licitaciones públicas de obras de agua y saneamiento convocadas por el Estado peruano, los econometristas analizan anomalías estadísticas en las ofertas de los postores buscando:}
\begin{opts}
    \item que todas las empresas postoras coticen exactamente al costo marginal unitario.
    \item varianza anormalmente baja entre ofertas y alta correlación de precios entre postores.
    \item un coeficiente $R^2$ igual a cero en la ecuación del presupuesto referencial.
    \item una distribución estrictamente bimodal en los precios de mercado regionales.
\end{opts}

\q{Un regulador implementa un modelo de Nowcasting para proyectar en tiempo real la recaudación tarifaria de las EPS durante episodios de crisis macroeconómica. Este método econométrico se caracteriza por:}
\begin{opts}
    \item asume que el término de perturbación estocástica futuro es conocido sin margen de incertidumbre.
    \item proyecta variables estructurales a 50 años aplicando exclusivamente el Teorema de Coase.
    \item utiliza únicamente modelos estáticos de equilibrio general computable sin rezagos dinámicos.
    \item estima el estado contemporáneo usando datos de alta frecuencia previos al reporte contable.
\end{opts}

\q{En las auditorías regulatorias de calidad bacteriológica del agua potable realizadas por SUNASS y DIGESA, la determinación del tamaño mínimo de muestra de grifos fiscalizados ($n$) para contrastar cumplimiento normativo exige balancear:}
\begin{opts}
    \item la depreciación contable lineal de los equipos de laboratorio y los costos fijos administrativos.
    \item la tasa de descuento intertemporal del BCRP y la volatilidad del tipo de cambio spot.
    \item el coeficiente $R^2$ ajustado de la regresión y el factor de inflación de varianza máximo.
    \item el error Tipo I (falsa sanción a la EPS) y el error Tipo II (no detectar agua contaminada).
\end{opts}


\section*{Bloque D. Ampliación: Teoría y Aplicaciones Prácticas (70\% / 30\%)}

\noindent\textit{Preguntas 91 a 180. Ampliación del banco oficial con 63 preguntas teóricas y conceptuales (91 a 153) y 27 aplicaciones cuantitativas prácticas (154 a 180).}

\subsection*{Parte I. Preguntas teóricas y conceptuales (91 a 153)}

\q{En el modelo lineal general $Y = X\beta + u$, la matriz generadora de residuos (o matriz aniquiladora) $M = I_n - X(X'X)^{-1}X'$ cumple simultáneamente las siguientes propiedades algebraicas:}
\begin{opts}
    \item Es simétrica, idempotente ($M^2 = M$), satisface $MX = 0$ y su traza es igual a $n - k$.
    \item Es diagonal, invertible y posee determinante estrictamente positivo igual a $n - k$ para cualquier muestra.
    \item Es ortogonal, definida positiva, satisface $MX = X$ y su rango es igual exactamente a $k$ parámetros.
    \item Es antisimétrica, singular, satisface $MY = X\hat{\beta}$ y su traza equivale exactamente a $k$ regresores.
\end{opts}

\q{De acuerdo con el Teorema de Frisch-Waugh-Lovell (FWL) aplicado a la regresión particionada $Y = X_1\beta_1 + X_2\beta_2 + u$, el estimador MCO $\hat{\beta}_1$ es algebraicamente idéntico al que se obtiene al:}
\begin{opts}
    \item Regresar $Y$ directamente sobre $X_1$ omitiendo $X_2$ y luego restar la covarianza muestral entre ambos bloques.
    \item Regresar los residuos de $Y$ sobre $X_2$ contra los residuos de $X_1$ sobre $X_2$ ($M_2 Y$ sobre $M_2 X_1$).
    \item Regresar $Y$ sobre los valores ajustados $\hat{X}_1$ obtenidos de proyectar $X_1$ sobre el subespacio de $X_2$.
    \item Estimar por separado $Y$ sobre $X_1$ y $Y$ sobre $X_2$ y ponderar ambos coeficientes por sus respectivas varianzas.
\end{opts}

\q{Cuando se estima un modelo de regresión lineal por MCO restringiendo el intercepto a cero (regresión a través del origen), una consecuencia algebraica directa es que:}
\begin{opts}
    \item El estimador de la pendiente mantiene exactamente la misma varianza que en el modelo con intercepto libre.
    \item Los estimadores MCO pasan a ser automáticamente consistentes incluso en presencia de endogeneidad severa.
    \item La suma de los residuos MCO ya no es necesariamente cero y el $R^2$ centrado puede resultar negativo.
    \item La descomposición $SCT = SCE + SCR$ se conserva intacta y el $R^2$ siempre permanece estrictamente entre $0$ y $1$.
\end{opts}

\q{Respecto a la relación entre el supuesto de exogeneidad estricta $E[u|X] = 0$ y la condición de ortogonalidad $E[X'u] = 0$ en el modelo lineal, es correcto afirmar que:}
\begin{opts}
    \item La condición $E[X'u] = 0$ es más restrictiva que $E[u|X] = 0$ y asegura insesgadez exacta en cualquier muestra finita.
    \item Ambos supuestos son estrictamente equivalentes y requieren necesariamente que los errores sigan una distribución normal exacta.
    \item Ninguno de los dos supuestos resulta necesario para la consistencia si se cumple la homocedasticidad condicional de los errores.
    \item El supuesto $E[u|X] = 0$ garantiza insesgadez en muestras finitas, mientras que $E[X'u] = 0$ basta para consistencia asintótica.
\end{opts}

\q{Si a los supuestos clásicos de Gauss-Markov se añade el supuesto de normalidad de las perturbaciones $u|X \sim N(0, \sigma^2 I_n)$, el estimador MCO $\hat{\beta}$ alcanza una propiedad de eficiencia más fuerte porque:}
\begin{opts}
    \item Se vuelve el estimador de mínima varianza únicamente dentro de la subclase de estimadores lineales e insesgados.
    \item Coincide con el estimador de Máxima Verosimilitud y alcanza la cota inferior de Cramér-Rao (es MIE / BUE).
    \item Elimina por completo el error de muestreo en muestras pequeñas sin necesidad de que exista rango completo en la matriz $X$.
    \item Permite que el estimador de la varianza residual de Máxima Verosimilitud $\tilde{\sigma}^2 = SCR/n$ sea estrictamente insesgado.
\end{opts}

\q{En una regresión lineal múltiple estimada por MCO, si se multiplican todos los valores del regresor $X_2$ por una constante $c = 1000$ (por ejemplo, pasando de miles de soles a soles) sin alterar $Y$ ni los demás regresores, entonces:}
\begin{opts}
    \item El coeficiente $\hat{\beta}_2$ no cambia, pero su error estándar se divide entre $1000$ y el $R^2$ se mantiene completamente inalterado.
    \item El coeficiente $\hat{\beta}_2$ se multiplica por $1000$, mientras que el estadístico $t_2$ y el $R^2$ aumentan proporcionalmente.
    \item Tanto $\hat{\beta}_2$ como su error estándar se dividen entre $1000$, dejando invariantes el estadístico $t_2$ y el $R^2$.
    \item Tanto $\hat{\beta}_2$ como el intercepto $\hat{\beta}_0$ se dividen entre $1000$, reduciendo el estadístico $F$ global del modelo.
\end{opts}

\q{¿Cuál es la diferencia conceptual y estadística fundamental entre la perturbación aleatoria $u_i$ y el residuo MCO $\hat{u}_i$ en una regresión lineal?}
\begin{opts}
    \item La perturbación $u_i$ es observable a partir de los datos muestrales, mientras que $\hat{u}_i$ es una variable latente teórica inobservable.
    \item Ambos coinciden idénticamente en muestras finitas siempre que el modelo incluya un intercepto y no exista multicolinealidad perfecta.
    \item Los residuos $\hat{u}_i$ siempre son independientes entre sí y homocedásticos, mientras que las perturbaciones $u_i$ están correlacionadas entre sí.
    \item La perturbación $u_i$ es inobservable con varianza $\sigma^2$; $\hat{u}_i = Y_i - \hat{Y}_i$ es muestral con $\text{Var}(\hat{u}|X) = \sigma^2 M$.
\end{opts}

\q{Al añadir una nueva variable explicativa $X_{k+1}$ a un modelo de regresión lineal múltiple estimado por MCO, el coeficiente de determinación ajustado $\bar{R}^2$ aumentará si y solo si:}
\begin{opts}
    \item El valor absoluto del estadístico $t$ asociado a la nueva variable incorporada es mayor que $1$ ($|t| > 1$).
    \item La correlación simple entre $Y$ y $X_{k+1}$ supera el valor del $R^2$ del modelo original sin restringir.
    \item El coeficiente estimado $\hat{\beta}_{k+1}$ resulta estrictamente positivo independientemente de su error estándar muestral.
    \item La nueva variable reduce la suma de cuadrados totales ($SCT$) manteniendo inalterados los grados de libertad residuales.
\end{opts}

\q{En un modelo semilogarítmico $\ln(Y_i) = \beta_0 + \beta_1 X_i + \delta D_i + u_i$, donde $D_i \in \{0, 1\}$ es una variable binaria (por ejemplo, contar con micromedidor de agua), el cambio porcentual exacto en $Y$ al pasar de $D_i = 0$ a $D_i = 1$ (ceteris paribus) es:}
\begin{opts}
    \item Exactamente $100 \cdot \delta\%$ para cualquier magnitud positiva o negativa que tome el parámetro $\delta$ en la población.
    \item Exactamente $100 \cdot \ln(1 + \delta)\%$, válido exclusivamente si el error sigue una distribución logística simétrica.
    \item Exactamente $100 \cdot (e^{\delta} - 1)\%$, aproximado por $100 \cdot \delta\%$ solo cuando $|\delta|$ es pequeño.
    \item Exactamente $100 \cdot (\delta / \bar{Y})\%$, evaluado en el promedio muestral de la variable dependiente transformada.
\end{opts}

\q{En una especificación cuadrática $Y_i = \beta_0 + \beta_1 X_i + \beta_2 X_i^2 + u_i$ con $\beta_1 > 0$ y $\beta_2 < 0$, el efecto marginal de $X$ sobre $Y$ y el valor de $X$ donde $E[Y|X]$ alcanza su máximo son, respectivamente:}
\begin{opts}
    \item Constante igual a $\beta_1 + 2\beta_2$ y un máximo ubicado en $X^* = -\beta_2 / (2\beta_1)$ sobre el eje horizontal.
    \item Igual a $\beta_1 + \beta_2 X$ y un punto de inflexión ubicado en $X^* = -\beta_1 / \beta_2$ para todo $X > 0$.
    \item Igual a $2\beta_1 X + \beta_2$ y un mínimo global ubicado en $X^* = \beta_1 / (2\beta_2)$ en la muestra.
    \item Igual a $\beta_1 + 2\beta_2 X$ y un máximo ubicado en $X^* = -\beta_1 / (2\beta_2)$.
\end{opts}

\q{Al contrastar una única restricción lineal ($q = 1$) sobre los parámetros de un modelo clásico con errores normales, la relación exacta entre el estadístico $F$ y el estadístico $t$ bilateral es:}
\begin{opts}
    \item $F_{1, n-k} = t_{n-k}^2$, por lo que ambas pruebas producen el mismo p-valor y decisión.
    \item $F_{1, n-k} = \sqrt{|t_{n-k}|}$, de modo que la prueba $F$ resulta siempre más conservadora que la prueba $t$.
    \item $F_{1, n-k} = (n - k) \cdot t_{n-k}^2$, razón por la cual solo coinciden en el límite asintótico de muestras grandes.
    \item No existe equivalencia exacta porque la prueba $F$ es unilateral a la izquierda y la prueba $t$ es siempre bilateral.
\end{opts}

\q{En econometría aplicada con bases de datos masivas (por ejemplo, millones de registros de facturación de agua), ¿por qué un coeficiente puede ser altamente significativo estadísticamente ($p < 0.001$) y, sin embargo, carecer de relevancia económica?}
\begin{opts}
    \item Porque al crecer $n$, el estimador MCO se vuelve sesgado y converge hacia un valor distinto del parámetro poblacional real.
    \item Porque al crecer $n$, los errores estándar tienden a cero y rechazan $H_0: \beta_j = 0$ incluso ante efectos minúsculos.
    \item Porque la distribución $t$ de Student deja de converger a la normal estándar cuando $n$ supera las mil observaciones.
    \item Porque el coeficiente de determinación $R^2$ tiende obligatoriamente a uno conforme aumenta el tamaño muestral del estudio.
\end{opts}

\q{La prueba de Chow para detectar un cambio estructural conocido en una fecha $t_0$ entre dos subperiodos equivale a realizar una prueba $F$ conjunta sobre:}
\begin{opts}
    \item La igualdad exclusiva de las varianzas residuales de ambos subperiodos sin restringir interceptos ni pendientes del modelo.
    \item La significancia de una única variable dummy de periodo que solo permite variar el intercepto entre subperiodos.
    \item La ausencia de autocorrelación de primer orden en los residuos agrupados del periodo completo de estimación econométrica.
    \item La significancia conjunta de una dummy de periodo y todas sus interacciones con los regresores.
\end{opts}

\q{En un modelo lineal con errores normales estimado por Máxima Verosimilitud, los tres contrastes asintóticos clásicos para $q$ restricciones lineales —Wald ($W$), Razón de Verosimilitud ($LR$) y Multiplicador de Lagrange ($LM$)— satisfacen:}
\begin{opts}
    \item La desigualdad $W \ge LR \ge LM$ en muestras finitas y convergen todos a $\chi^2_q$ bajo $H_0$.
    \item La desigualdad $LM \ge LR \ge W$ en muestras finitas y convergen a distintas distribuciones $t$ de Student asintóticamente.
    \item La igualdad exacta $W = LR = LM$ para cualquier tamaño de muestra finita siempre que existan $q \ge 2$ restricciones.
    \item Que únicamente $LM$ requiere estimar el modelo no restringido, mientras que $W$ utiliza exclusivamente el modelo restringido.
\end{opts}

\q{Si un investigador incluye en una regresión lineal una variable explicativa completamente irrelevante (cuyo verdadero parámetro poblacional es $\beta_{irr} = 0$) que está correlacionada con los demás regresores, el estimador MCO de los parámetros relevantes:}
\begin{opts}
    \item Se vuelve sesgado e inconsistente debido a la pérdida de grados de libertad en la estimación de la matriz $X'X$.
    \item Permanece insesgado y consistente, pero pierde eficiencia al aumentar las varianzas de los estimadores.
    \item Aumenta su eficiencia asintótica porque el coeficiente de determinación $R^2$ del modelo no disminuye al incorporar regresores.
    \item Deja de ser lineal en $Y$ y viola el supuesto de exogeneidad estricta requerido por el teorema clásico de Gauss-Markov.
\end{opts}

\q{Suponga que el modelo verdadero es $Y = \beta_0 + \beta_1 X_1 + \beta_2 X_2 + u$ con $\beta_2 \neq 0$, pero se omite $X_2$ al estimar por MCO. Si $X_1$ y $X_2$ son estrictamente ortogonales en la muestra ($\text{Cov}(X_1, X_2) = 0$), entonces:}
\begin{opts}
    \item $\hat{\beta}_1$ sigue siendo sesgado, aunque el estimador de la varianza residual $s^2$ permanece completamente insesgado.
    \item Tanto $\hat{\beta}_1$ como el estimador de la varianza residual $s^2$ se mantienen insesgados y asintóticamente eficientes.
    \item $\hat{\beta}_1$ permanece insesgado, pero el estimador de la varianza residual $s^2$ queda sesgado al alza.
    \item $\hat{\beta}_1$ resulta inconsistente y la prueba $t$ asociada tiende a rechazar $H_0$ con mayor frecuencia de la real.
\end{opts}

\q{Al comparar los criterios de información de Akaike (AIC) y Bayesiano de Schwarz (BIC) para seleccionar entre modelos no anidados con $n \ge 8$ observaciones, es correcto señalar que:}
\begin{opts}
    \item BIC impone una penalización mayor por parámetro ($\ln n > 2$) y es consistente para seleccionar el modelo verdadero.
    \item AIC penaliza más fuertemente la inclusión de parámetros adicionales y tiende siempre a elegir modelos mucho más parsimoniosos.
    \item Ambos criterios seleccionan el modelo con el valor numérico más alto posible de la función objetivo evaluada en los residuos.
    \item El criterio BIC ignora por completo el tamaño muestral $n$ y depende exclusivamente del número de regresores $k$ incluidos.
\end{opts}

\q{En una prueba de hipótesis econométrica para evaluar si la calidad del agua potable cumple el estándar regulatorio, mantener fijo el tamaño muestral $n$ y reducir el nivel de significancia $\alpha$ de $5\%$ a $1\%$ provoca:}
\begin{opts}
    \item Una disminución simultánea de la probabilidad de cometer Error Tipo I y de la probabilidad de cometer Error Tipo II.
    \item Una menor probabilidad de Error Tipo I, pero mayor probabilidad de Error Tipo II (menor potencia).
    \item Un incremento en la potencia estadística $(1 - \beta)$ sin alterar el tamaño de la región crítica del contraste bilateral.
    \item Una reducción del error estándar del estimador y un estrechamiento automático del intervalo de confianza correspondiente.
\end{opts}

\q{En presencia de multicolinealidad imperfecta severa entre dos regresores de un modelo de costos de una EPS, un síntoma clásico en la salida econométrica consiste en observar:}
\begin{opts}
    \item Un $R^2$ muy bajo y un estadístico $F$ global no significativo, acompañados de estadísticos $t$ individuales muy elevados.
    \item Estimadores MCO severamente sesgados e inconsistentes cuyo signo se invierte sistemáticamente al aumentar el tamaño de muestra $n$.
    \item Un $R^2$ alto y un estadístico $F$ global significativo, junto con estadísticos $t$ individuales bajos.
    \item Residuos altamente correlacionados con las variables explicativas, violando la ortogonalidad algebraica $X'\hat{u} = 0$ del estimador.
\end{opts}

\q{En la expresión de la varianza del estimador MCO $\text{Var}(\hat{\beta}_j | X) = \frac{\sigma^2}{SCT_j (1 - R_j^2)}$, el término $R_j^2$ representa:}
\begin{opts}
    \item El coeficiente de determinación de la regresión principal de $Y$ sobre todos los regresores del modelo estimado.
    \item La proporción de la varianza residual explicada por los términos cuadráticos en la prueba general de heterocedasticidad.
    \item El cuadrado de la correlación parcial entre la variable dependiente $Y$ y la variable explicativa $X_j$ seleccionada.
    \item El coeficiente de determinación de la regresión auxiliar de $X_j$ sobre los demás regresores.
\end{opts}

\q{Una cautela importante al interpretar un rechazo de la hipótesis nula en la prueba general de White es que este estadístico puede resultar significativo no solo por heterocedasticidad pura, sino también por:}
\begin{opts}
    \item La presencia de multicolinealidad imperfecta moderada entre dos variables explicativas continuas del modelo lineal.
    \item Errores de especificación en la forma funcional de la media condicional (como omisión de no linealidades).
    \item Un cambio de unidades de medida en la variable dependiente $Y$ aplicado previamente a la estimación por MCO.
    \item La falta de normalidad asintótica de las variables explicativas cuando el tamaño muestral es suficientemente elevado.
\end{opts}

\q{¿En qué consiste el procedimiento de la prueba de Goldfeld-Quandt para contrastar heterocedasticidad vinculada a una variable $X_j$?}
\begin{opts}
    \item Dividir la muestra aleatoriamente en dos mitades iguales y aplicar una prueba de Kolmogorov-Smirnov sobre los residuos.
    \item Regresar los residuos al cuadrado sobre todas las variables explicativas, sus cuadrados y todos sus productos cruzados.
    \item Ordenar según $X_j$, omitir $c$ datos centrales, estimar dos regresiones extremas y calcular $F = SCR_2 / SCR_1$.
    \item Estimar una regresión auxiliar del valor absoluto de los residuos sobre el logaritmo natural de los valores predichos $\hat{Y}_i$.
\end{opts}

\q{A diferencia de Mínimos Cuadrados Generalizados (MCG) con matriz $\Omega$ conocida, el estimador de Mínimos Cuadrados Generalizados Factibles (MCGF / FGLS) que estima los parámetros de la varianza desconocida:}
\begin{opts}
    \item Es siempre estrictamente insesgado en muestras finitas y exactamente igual al estimador MCO convencional sin ponderar.
    \item Solo puede aplicarse en modelos de series de tiempo macroeconómicas y carece de validez en datos de corte transversal.
    \item Requiere obligatoriamente conocer a priori la varianza poblacional exacta $\sigma_i^2$ de cada observación individual de la muestra.
    \item No es necesariamente insesgado en muestras pequeñas, pero es consistente y asintóticamente más eficiente que MCO.
\end{opts}

\q{¿Qué distingue conceptualmente a un proceso ARCH/GARCH respecto de la heterocedasticidad clásica de corte transversal?}
\begin{opts}
    \item En ARCH/GARCH la varianza condicional $\sigma_t^2$ evoluciona según shocks pasados ($\varepsilon_{t-1}^2$) y varianzas rezagadas.
    \item En ARCH/GARCH la varianza incondicional siempre diverge al infinito, mientras que la varianza condicional permanece constante en el tiempo.
    \item Los modelos ARCH/GARCH corrigen únicamente la endogeneidad por variables omitidas en series macroeconómicas trimestrales.
    \item En los modelos ARCH/GARCH el término de error está correlacionado serialmente en media ($E[\varepsilon_t \varepsilon_{t-1}] \neq 0$).
\end{opts}

\q{Si los errores de una regresión de series de tiempo exhiben autocorrelación serial negativa ($\rho < 0$) en un proceso $AR(1)$, el estadístico de Durbin-Watson ($DW$) se ubicará aproximadamente en el intervalo:}
\begin{opts}
    \item Entre $0$ y $2$, tendiendo hacia $0$ conforme el coeficiente $\rho$ se acerca al límite $-1$.
    \item Exactamente en $2$, independientemente del signo y de la magnitud del coeficiente autorregresivo $\rho$.
    \item Entre $2$ y $4$, tendiendo a $4$ conforme $\rho$ se acerca a $-1$.
    \item En valores estrictamente negativos, dado que el estadístico cumple aproximadamente $DW \approx -2(1 - \hat{\rho})$.
\end{opts}

\q{Cuando un modelo dinámico incluye la variable dependiente rezagada $Y_{t-1}$ como regresor, el estadístico $d$ de Durbin-Watson está sesgado hacia $2$. Para contrastar autocorrelación $AR(1)$ en este contexto, una alternativa válida en muestras grandes es:}
\begin{opts}
    \item La prueba de normalidad de Jarque-Bera evaluada en la regresión auxiliar en primeras diferencias.
    \item La prueba de Goldfeld-Quandt aplicada sobre los residuos ordenados de forma estrictamente cronológica.
    \item El factor de inflación de varianza (VIF) calculado específicamente sobre la variable dependiente rezagada.
    \item El estadístico $h$ de Durbin o, de forma más general, la prueba LM de Breusch-Godfrey.
\end{opts}

\q{Al corregir autocorrelación $AR(1)$ mediante MCG factibles, la ventaja principal del procedimiento de Prais-Winsten frente al algoritmo iterativo de Cochrane-Orcutt radica en que:}
\begin{opts}
    \item Prais-Winsten conserva la primera observación ($t=1$) escalándola por $\sqrt{1 - \hat{\rho}^2}$, mejorando eficiencia.
    \item Cochrane-Orcutt exige conocer el valor poblacional exacto de $\rho$, mientras que Prais-Winsten no requiere estimar ningún parámetro.
    \item Prais-Winsten corrige automáticamente problemas de endogeneidad y simultaneidad que el método de Cochrane-Orcutt deja intactos.
    \item Cochrane-Orcutt solo funciona cuando $\rho < 0$, en tanto Prais-Winsten es exclusivo para procesos con raíz unitaria.
\end{opts}

\q{Al utilizar el estimador de varianzas HAC de Newey-West en una regresión de series de tiempo, los coeficientes puntuales $\hat{\beta}$ obtenidos respecto de MCO tradicional:}
\begin{opts}
    \item Cambian sustancialmente porque Newey-West aplica una transformación cuasi-diferenciada a todas las variables del modelo.
    \item Permanecen idénticos a los de MCO; únicamente se corrigen la matriz de covarianzas y los errores estándar.
    \item Se vuelven automáticamente consistentes incluso cuando existe simultaneidad severa entre los regresores y el término de error.
    \item Alcanzan la eficiencia exacta de Mínimos Cuadrados Generalizados en muestras finitas pequeñas sin estimar $\rho$.
\end{opts}

\q{En el modelo de regresión lineal simple $Y_i = \beta_0 + \beta_1 X_i^* + u_i$, contraste el efecto de un error de medición clásico (media cero e independiente del verdadero valor) en $X_i^*$ frente a un error de medición clásico en $Y_i$:}
\begin{opts}
    \item El error en $Y_i$ sesga $\hat{\beta}_1$ hacia cero de manera inconsistente, mientras que el error en $X_i^*$ solo incrementa la varianza residual.
    \item Ninguno de los dos errores afecta la consistencia asintótica de MCO siempre que el tamaño muestral $n$ sea suficientemente grande.
    \item Ambos tipos de error de medición generan exactamente el mismo sesgo de atenuación proporcional a la varianza del ruido aleatorio.
    \item El error en $X_i^*$ genera sesgo de atenuación hacia cero; el error en $Y_i$ mantiene la insesgadez pero infla varianzas.
\end{opts}

\q{En un sistema lineal de ecuaciones simultáneas (por ejemplo, oferta y demanda de agua), respecto de las condiciones de identificación de una ecuación estructural es correcto afirmar que:}
\begin{opts}
    \item La condición de orden (exógenas excluidas $\ge$ endógenas incluidas menos una) es necesaria; la de rango es necesaria y suficiente.
    \item La condición de orden es suficiente por sí sola para garantizar que todos los parámetros estructurales estén exactamente identificados.
    \item Una ecuación exactamente identificada no puede estimarse mediante MC2E y requiere obligatoriamente aplicar Mínimos Cuadrados Ordinarios.
    \item La condición de rango exige que todas las variables exógenas del sistema aparezcan incluidas simultáneamente en cada ecuación estructural.
\end{opts}

\q{Bajo efectos de tratamiento heterogéneos y un instrumento binario válido que satisface el supuesto de monotonicidad (ausencia de desafiantes o *defiers*), el estimador de Variables Instrumentales (estimador de Wald) identifica:}
\begin{opts}
    \item El Efecto Promedio del Tratamiento (ATE) sobre toda la población objetivo independientemente de su respuesta al instrumento asignado.
    \item El Efecto Local Promedio del Tratamiento (LATE) para la subpoblación de cumplidores (*compliers*).
    \item El Efecto Promedio sobre los Nunca Tratados (*never-takers*) evaluado exactamente en el umbral de elegibilidad.
    \item Una cota inferior del sesgo de selección muestral en observables sin ninguna interpretación causal directa.
\end{opts}

\q{Cuando se emplean instrumentos débiles (por ejemplo, un estadístico $F$ de la primera etapa muy inferior a $10$) en una estimación por Mínimos Cuadrados en Dos Etapas (MC2E), la principal consecuencia econométrica es que:}
\begin{opts}
    \item El estimador MC2E converge más rápidamente al parámetro verdadero y presenta errores estándar mucho más precisos que MCO.
    \item El estimador MC2E se vuelve exactamente igual al estimador de Efectos Fijos y elimina por completo el sesgo de variable omitida.
    \item MC2E exhibe un fuerte sesgo en muestras finitas hacia MCO inconsistente y pierde cobertura en sus intervalos.
    \item La prueba de sobreidentificación de Sargan rechaza automáticamente la hipótesis nula con probabilidad exactamente igual a uno.
\end{opts}

\q{¿Por qué no es posible contrastar estadísticamente la condición de exogeneidad del instrumento (restricción de exclusión $\text{Cov}(Z, u) = 0$) en un modelo exactamente identificado (un instrumento para una variable endógena)?}
\begin{opts}
    \item Porque $Z'\hat{u}_{VI} = 0$ se cumple por construcción, sin grados de libertad para probar sobreidentificación.
    \item Porque la condición de relevancia $\text{Cov}(Z, X) \neq 0$ impide calcular los residuos estructurales de la segunda etapa.
    \item Porque el estimador de primera etapa es inconsistente cuando el número de instrumentos iguala al número de regresores endógenos.
    \item Porque la distribución asintótica del estimador VI en modelos exactamente identificados no depende de la covarianza con el error.
\end{opts}

\q{En un modelo lineal sobreidentificado ($L > m$ instrumentos) con perturbaciones heterocedásticas, ¿por qué el estimador óptimo del Método Generalizado de Momentos (GMM en dos etapas) supera asintóticamente a MC2E?}
\begin{opts}
    \item Porque GMM no requiere que los instrumentos satisfagan la condición de relevancia estadística en la primera etapa del modelo.
    \item Porque GMM usa la matriz óptima de pesos de momentos, logrando menor varianza asintótica que MC2E.
    \item Porque MC2E resulta inconsistente en presencia de heterocedasticidad, mientras que GMM corrige el sesgo asintótico.
    \item Porque GMM transforma automáticamente la variable dependiente en logaritmos y elimina cualquier correlación entre los instrumentos.
\end{opts}

\q{Al evaluar la validez conjunta de las restricciones de sobreidentificación en un modelo con heterocedasticidad, la ventaja del estadístico $J$ de Hansen frente al estadístico clásico de Sargan es que:}
\begin{opts}
    \item El estadístico de Sargan contrasta relevancia de los instrumentos en primera etapa, mientras que Hansen-J contrasta normalidad.
    \item El estadístico $J$ de Hansen permite contrastar la exogeneidad incluso en modelos exactamente identificados ($L = m$).
    \item Sargan requiere homocedasticidad para seguir $\chi^2_{L-m}$, mientras que Hansen-J es robusto a heterocedasticidad.
    \item El estadístico $J$ de Hansen no pierde potencia estadística aunque se incluyan decenas de instrumentos débiles en la estimación.
\end{opts}

\q{En inferencia causal aplicada, incluir como covariable de control una variable que a su vez es un resultado afectado por el tratamiento (un 'mal control' o *bad control*) produce:}
\begin{opts}
    \item Una reducción garantizada del error estándar sin afectar en absoluto la insesgadez del efecto causal estimado.
    \item Un aumento de la validez externa del estimador al restringir la muestra exclusivamente a los individuos tratados.
    \item Una solución exacta al problema de endogeneidad causado por variables omitidas estrictamente invariantes en el tiempo.
    \item Sesgo de selección post-tratamiento al bloquear parte del mecanismo causal o abrir rutas espurias de colisionador.
\end{opts}

\q{En un modelo de datos de panel de empresas prestadoras de saneamiento (EPS) estimado por Efectos Fijos mediante la transformación *within* ($y_{it} - \bar{y}_i$), ¿por qué no es posible estimar el coeficiente de una variable estrictamente invariante en el tiempo (como la altitud de la ciudad)?}
\begin{opts}
    \item Porque las variables invariantes en el tiempo siempre presentan autocorrelación serial perfecta de segundo orden en paneles.
    \item Porque al restar la media temporal de cada EPS, toda variable constante en el tiempo se anula ($x_i - \bar{x}_i = 0$).
    \item Porque el estimador *within* utiliza exclusivamente la variación entre distintas EPS y descarta la variación temporal.
    \item Porque la presencia de variables invariantes en el tiempo viola la condición de homocedasticidad en datos de panel.
\end{opts}

\q{En un panel lineal sin variables dependientes rezagadas, la relación entre el estimador *within* (variables centradas en su media individual) y el estimador MCO con variables ficticias individuales (LSDV) es que:}
\begin{opts}
    \item El estimador *within* produce coeficientes más pequeños porque omite por completo la variación del intercepto común.
    \item LSDV es consistente únicamente bajo el supuesto de efectos aleatorios, mientras que *within* es consistente bajo efectos fijos.
    \item Ambos producen idénticos coeficientes $\hat{\beta}_{FE}$, aunque LSDV consume explícitamente $N$ grados de libertad.
    \item LSDV solo coincide numéricamente con *within* cuando el número de periodos $T$ tiende a infinito y $N$ permanece fijo.
\end{opts}

\q{Al comparar el estimador de Efectos Fijos (FE) y el estimador de Primeras Diferencias (FD) en un panel con $T > 2$ periodos, ¿en qué escenario resulta más eficiente el estimador FD que el estimador FE?}
\begin{opts}
    \item Cuando los errores idiosincráticos $u_{it}$ no presentan autocorrelación serial de ningún orden y son homocedásticos.
    \item Cuando el efecto inobservable $\alpha_i$ no está correlacionado con ninguna de las variables explicativas $X_{it}$ incluidas.
    \item Cuando las variables explicativas son estrictamente constantes a lo largo de todos los periodos temporales del panel.
    \item Cuando $u_{it}$ sigue una caminata aleatoria ($u_{it} = u_{it-1} + e_{it}$), haciendo que $\Delta u_{it}$ sea ruido blanco.
\end{opts}

\q{El estimador de Efectos Aleatorios (RE) en datos de panel se obtiene aplicando MCG mediante una transformación de cuasi-centrado ($y_{it} - \theta \bar{y}_i$). Cuando la varianza del efecto individual $\sigma_{\alpha}^2$ es muy grande en relación con $\sigma_u^2$ (o $T \to \infty$), el factor $\theta$:}
\begin{opts}
    \item Tiende a $1$, haciendo que Efectos Aleatorios converja al estimador de Efectos Fijos (*within*).
    \item Tiende a $0$, haciendo que el estimador de Efectos Aleatorios converja al estimador de MCO agrupado (*pooled OLS*).
    \item Tiende a infinito, invalidando la existencia de la matriz de varianzas-covarianzas de los errores compuestos.
    \item Permanece fijo exactamente en $0.5$, ponderando por igual la variación *within* y la variación *between* del panel.
\end{opts}

\q{En estudios empíricos con datos de panel de EPS (por ejemplo, $50$ empresas observadas durante $15$ años), ¿por qué se recomienda calcular errores estándar agrupados (*clustered*) a nivel de EPS tras estimar Efectos Fijos?}
\begin{opts}
    \item Para corregir el sesgo en los coeficientes $\hat{\beta}_{FE}$ causado por la presencia de variables omitidas cambiantes en el tiempo.
    \item Para reducir artificialmente los errores estándar y asegurar que los coeficientes resulten estadísticamente significativos al $5\%$.
    \item Para permitir heterocedasticidad y correlación serial arbitraria dentro de los residuos de cada EPS.
    \item Para transformar el estimador de Efectos Fijos en un estimador de Máxima Verosimilitud con distribución normal exacta.
\end{opts}

\q{En un modelo de panel dinámico corto ($N$ grande, $T$ pequeño) que incluye la variable dependiente rezagada $y_{it-1}$ junto con efectos fijos individuales $\alpha_i$, estimar mediante el estimador *within* tradicional produce:}
\begin{opts}
    \item Un estimador completamente insesgado y consistente debido a que la transformación *within* elimina el efecto fijo $\alpha_i$.
    \item Estimadores consistentes en muestras finitas únicamente si el coeficiente autorregresivo verdadero es estrictamente igual a cero.
    \item Un sesgo que desaparece automáticamente al reemplazar los errores estándar clásicos por errores agrupados por individuo.
    \item El sesgo de Nickell por correlación entre $\tilde{y}_{it-1}$ y $\tilde{u}_{it}$; suele resolverse con GMM de Arellano-Bond.
\end{opts}

\q{En un panel desbalanceado de hogares evaluados por SUNASS donde algunas familias abandonan la muestra en periodos posteriores (atrición), el estimador de Efectos Fijos conserva la consistencia siempre que:}
\begin{opts}
    \item La atrición dependa solo de observables incluidos o de $\alpha_i$, y no del error idiosincrático $u_{it}$.
    \item La probabilidad de permanecer en la muestra sea estrictamente idéntica a $1$ para todos los hogares encuestados inicialmente.
    \item Se eliminen todos los periodos anteriores de los hogares que abandonan para forzar mecánicamente un panel balanceado.
    \item La atrición esté fuertemente correlacionada con choques transitorios no observados en el consumo mensual de agua $u_{it}$.
\end{opts}

\q{¿Cuál es la diferencia fundamental entre una serie estacionaria en tendencia (TS) y una serie estacionaria en diferencias (DS, con raíz unitaria) frente a un choque aleatorio transitorio $\varepsilon_t$?}
\begin{opts}
    \item En una serie TS el choque tiene un efecto permanente sobre la trayectoria futura, mientras que en una serie DS se disipa rápidamente.
    \item En una serie TS el choque es transitorio (revierte a la tendencia); en una serie DS su efecto es permanente.
    \item Ambas series poseen varianza incondicional constante en el tiempo y media completamente independiente del periodo $t$.
    \item Una serie TS requiere diferenciarse dos veces para volverse estacionaria, en tanto una DS solo requiere restar una tendencia lineal.
\end{opts}

\q{Para un proceso autorregresivo $AR(p)$ representado como $\phi(L) y_t = c + \varepsilon_t$, la condición necesaria y suficiente para que el proceso sea estacionario en covarianza es que:}
\begin{opts}
    \item Al menos una raíz del polinomio $\phi(z) = 0$ sea igual a $1$ y las demás sean números reales negativos.
    \item Todas las raíces del polinomio característico $\phi(z) = 0$ caigan estrictamente dentro del círculo unitario ($|z| < 1$).
    \item La suma de los coeficientes autorregresivos sea estrictamente igual a la unidad ($\sum_{i=1}^p \phi_i = 1$).
    \item Todas las raíces del polinomio característico $\phi(z) = 0$ caigan fuera del círculo unitario ($|z| > 1$).
\end{opts}

\q{Al identificar un proceso estacionario $ARMA(p, q)$ mediante la metodología de Box-Jenkins, el patrón teórico característico de un proceso puro de medias móviles $MA(q)$ consiste en que:}
\begin{opts}
    \item La función de autocorrelación (FAC) se trunca después del rezago $q$, mientras la FACP decae gradualmente.
    \item La función de autocorrelación (FAC) decae exponencialmente y la autocorrelación parcial (FACP) se trunca exactamente después del rezago $q$.
    \item Tanto la FAC como la FACP se mantienen constantes e iguales a $1$ para todos los rezagos $k \ge 1$ del correlograma.
    \item Tanto la FAC como la FACP se anulan abruptamente a partir del primer rezago sin importar el valor del orden $q$.
\end{opts}

\q{En la prueba de raíz unitaria de Dickey-Fuller Aumentada (ADF), ¿con qué propósito se incorporan rezagos de la primera diferencia ($\Delta y_{t-1}, \dots, \Delta y_{t-p}$) en la ecuación auxiliar?}
\begin{opts}
    \item Para convertir la distribución asintótica del estadístico de prueba en una distribución $t$ de Student estándar simétrica.
    \item Para absorber la autocorrelación serial residual y lograr que el error de la ecuación auxiliar sea ruido blanco.
    \item Para invertir la hipótesis nula de modo que $H_0$ pase a representar estacionariedad en sentido débil de la serie.
    \item Para garantizar que el coeficiente de determinación $R^2$ de la regresión auxiliar supere al estadístico Durbin-Watson.
\end{opts}

\q{Según el Teorema de Representación de Granger, si dos series no estacionarias $y_t$ y $x_t$ son integradas de orden uno $I(1)$ y están cointegradas, entonces necesariamente:}
\begin{opts}
    \item Su relación de largo plazo es espuria y solo puede modelarse estimando un VAR exclusivamente en primeras diferencias.
    \item El estimador MCO en niveles resulta inconsistente y converge a cero conforme el tamaño muestral tiende a infinito.
    \item Admiten una representación mediante un Modelo de Corrección de Errores (MCE) que vincula corto y largo plazo.
    \item Ambas series carecen por completo de causalidad en el sentido de Granger en cualquiera de las dos direcciones posibles.
\end{opts}

\q{En el Modelo de Corrección de Errores $\Delta y_t = \gamma \Delta x_t + \alpha (y_{t-1} - \beta_0 - \beta_1 x_{t-1}) + \varepsilon_t$, ¿qué signo y rango debe tener el parámetro de ajuste $\alpha$ para que el sistema converja de forma estable hacia el equilibrio de largo plazo?}
\begin{opts}
    \item Debe ubicarse en $-1 < \alpha < 0$ (o en general $-2 < \alpha < 0$) para corregir el desequilibrio previo.
    \item Debe ser estrictamente positivo ($\alpha > 0$) para amplificar progresivamente las desviaciones respecto del equilibrio.
    \item Debe ser exactamente igual a cero ($\alpha = 0$), indicando que la serie no responde a los desequilibrios previos.
    \item Debe ser menor que $-2$ para garantizar que la varianza de $\Delta y_t$ converja a cero en un único periodo.
\end{opts}

\q{En el marco de un modelo de Vectores Autorregresivos (VAR) estacionario, afirmar que la variable $X_t$ 'causa en el sentido de Granger' a $Y_t$ significa formalmente que:}
\begin{opts}
    \item Existe una relación estructural causal demostrada que descarta la intervención de terceras variables omitidas en la economía.
    \item Los rezagos de $X_t$ mejoran significativamente el pronóstico de $Y_t$ más allá de la historia pasada de $Y_t$.
    \item La correlación contemporánea entre $X_t$ y $Y_t$ es estrictamente positiva y superior al $50\%$ en la muestra.
    \item Las series $X_t$ y $Y_t$ son obligatoriamente no estacionarias $I(1)$ y comparten una raíz unitaria común en el sistema.
\end{opts}

\q{Al calcular Funciones de Impulso-Respuesta (IRF) ortogonalizadas en un modelo VAR mediante la descomposición de Cholesky, una limitación metodológica bien conocida es que:}
\begin{opts}
    \item El método solo permite analizar choques de dos variables simultáneamente y no admite sistemas de tres o más ecuaciones.
    \item La descomposición de Cholesky impide que las variables del sistema VAR sean estacionarias en sentido débil.
    \item Las respuestas dependen del orden de las variables cuando los residuos reducidos tienen correlación contemporánea.
    \item Las funciones de impulso-respuesta obtenidas por Cholesky nunca convergen a cero aun cuando el VAR sea estable.
\end{opts}

\q{En los modelos de respuesta binaria Logit y Probit, dado que la probabilidad predicha $P(Y=1|X) = G(X'\beta)$ es no lineal en $X$, el efecto marginal de un regresor continuo $X_j$ sobre dicha probabilidad:}
\begin{opts}
    \item Es constante para todos los individuos de la muestra e idénticamente igual al coeficiente estimado $\hat{\beta}_j$.
    \item Se estima mediante Mínimos Cuadrados Ponderados en lugar de Máxima Verosimilitud para garantizar insesgadez exacta.
    \item Tiene siempre signo opuesto al del coeficiente estimado $\hat{\beta}_j$ cuando el índice $X'\beta$ toma valores negativos.
    \item Depende de $X$ mediante $g(X'\beta)\beta_j$, por lo que suele resumirse en la media (MEM) o en promedio (AME).
\end{opts}

\q{En un modelo Logit estimado para explicar la probabilidad de morosidad en el recibo de agua, la exponencial del coeficiente $e^{\hat{\beta}_j}$ se interpreta directamente como:}
\begin{opts}
    \item El incremento exacto en puntos porcentuales de la probabilidad de morosidad ante un aumento unitario de $X_j$.
    \item El factor multiplicativo en los *odds* ($\frac{p}{1-p}$) ante un incremento unitario de $X_j$.
    \item La elasticidad constante de la probabilidad de morosidad respecto de la variable explicativa $X_j$ en toda la muestra.
    \item El valor esperado del error latente logístico condicionado a que el usuario incurra efectivamente en morosidad.
\end{opts}

\q{En microeconometría aplicada (por ejemplo, modelar el gasto mensual de los hogares en agua embotellada, donde una fracción importante registra cero gasto), ¿qué diferencia existe entre una muestra censurada (modelo Tobit) y una muestra truncada?}
\begin{opts}
    \item La censura ocurre únicamente en series de tiempo financieras, mientras que el truncamiento es exclusivo de datos de panel.
    \item En una muestra truncada MCO produce estimadores insesgados, mientras que en una muestra censurada MCO es siempre inconsistente.
    \item En censura se observan las $X$ de toda la muestra; en truncamiento se pierden $Y$ y $X$ fuera del umbral.
    \item En una muestra censurada la variable dependiente es siempre binaria, en tanto en una truncada es siempre discreta de conteo.
\end{opts}

\q{Al modelar el número mensual de roturas de tuberías en redes de distribución de una EPS mediante una regresión de Poisson, si los datos presentan sobredispersión severa ($\text{Var}(Y|X) > E[Y|X]$), la principal consecuencia y alternativa estándar son:}
\begin{opts}
    \item El estimador MV de Poisson se vuelve inconsistente en media y debe reemplazarse obligatoriamente por un modelo Logit binario.
    \item Los errores estándar clásicos de Poisson quedan sobreestimados y conviene aplicar un modelo Tobit censurado en cero.
    \item La sobredispersión viola el supuesto de rango completo y hace imposible maximizar la función de log-verosimilitud de Poisson.
    \item La media sigue consistente (cuasi-MV), pero los errores estándar clásicos subestiman la varianza; se usa Binomial Negativa.
\end{opts}

\q{El pseudo-$R^2$ de McFadden en modelos de elección discreta estimados por Máxima Verosimilitud se define como $1 - \frac{\ln L_{ur}}{\ln L_0}$, donde $\ln L_0$ es la log-verosimilitud del modelo solo con intercepto. Respecto de este indicador es correcto afirmar que:}
\begin{opts}
    \item Toma valores en $[0, 1)$, acercándose a $0$ si los regresores no mejoran el ajuste y creciendo conforme $\ln L_{ur} \to 0$.
    \item Mide exactamente la proporción de la varianza de la variable dependiente explicada por los regresores igual que el $R^2$ de MCO.
    \item Puede tomar valores mayores que $1$ cuando el modelo predice correctamente más del $90\%$ de las observaciones binarias.
    \item Disminuye automáticamente cada vez que se añade una nueva variable explicativa relevante al modelo Logit o Probit.
\end{opts}

\q{En la evaluación de impacto mediante emparejamiento por puntaje de propensión (*Propensity Score Matching*, PSM), los dos supuestos fundamentales de identificación que permiten estimar el efecto causal promedio sobre los tratados (ATT) son:}
\begin{opts}
    \item Tendencias paralelas en el periodo pre-tratamiento y ausencia de correlación serial en los errores idiosincráticos.
    \item Existencia de un umbral continuo exógeno y continuidad de la función de densidad alrededor del punto de corte.
    \item Independencia condicional en observables (CIA) y existencia de soporte común (*common support*) entre grupos.
    \item Relevancia de un instrumento excluido en la primera etapa y normalidad exacta de la variable de resultado.
\end{opts}

\q{En un diseño de Regresión Discontinua Difusa (*Fuzzy RDD*), donde cruzar el umbral de puntaje SISFOH incrementa la probabilidad de recibir un subsidio de agua pero no de $0$ a $1$ de manera determinística, el efecto causal local en el umbral se estima mediante:}
\begin{opts}
    \item La diferencia simple de medias entre todos los hogares por debajo y por encima del puntaje de corte SISFOH.
    \item La exclusión de todos los hogares que no cumplieron la regla de asignación para convertir el diseño en un *Sharp RDD*.
    \item Un modelo de Efectos Aleatorios con errores estándar clásicos bajo el supuesto de homocedasticidad condicional.
    \item Variables Instrumentales: el cociente entre el salto en el resultado y el salto en la probabilidad de tratamiento.
\end{opts}

\q{¿Qué hipótesis evalúa la prueba de densidad de McCrary en un diseño de Regresión Discontinua (RDD) aplicado a la focalización de subsidios cruzados mediante el índice SISFOH?}
\begin{opts}
    \item Si hay discontinuidad en la densidad de la variable de asignación en el umbral (manipulación del puntaje).
    \item Si la variable de resultado $Y$ presenta heterocedasticidad condicional a ambos lados del punto de corte establecido.
    \item Si las tendencias temporales de los grupos tratado y de control eran paralelas antes de la reforma tarifaria.
    \item Si los residuos del polinomio local siguen una distribución normal estándar dentro de la ventana óptima de estimación.
\end{opts}

\q{En la literatura reciente de evaluación de impacto, cuando una política regulatoria (como la incorporación de EPS al Régimen de Apoyo Transitorio) se adopta de forma escalonada en distintos años (*staggered DiD*), estimar un modelo tradicional de Efectos Fijos de Dos Vías (TWFE) puede resultar sesgado si:}
\begin{opts}
    \item El efecto del tratamiento es estrictamente constante y homogéneo entre todas las cohortes y todos los periodos.
    \item Los efectos son heterogéneos entre cohortes o en el tiempo, pues TWFE usa ya tratados como controles.
    \item El número de periodos previos al tratamiento supera al número de periodos posteriores disponibles en el panel.
    \item Se utilizan errores estándar agrupados a nivel de empresa prestadora en lugar de errores robustos clásicos de White.
\end{opts}

\q{Al realizar benchmarking de eficiencia de costos entre las EPS peruanas para fijar metas de gestión, la diferencia metodológica esencial entre el Análisis Envolvente de Datos (DEA) y el Análisis de Frontera Estocástica (SFA) es que:}
\begin{opts}
    \item DEA es un método paramétrico que asume una función de costos Translog, mientras que SFA es determinístico y lineal.
    \item DEA únicamente puede medir eficiencia técnica con un solo insumo, mientras que SFA exige datos de series de tiempo puras.
    \item SFA no requiere especificar ninguna forma funcional ni distribución de errores, a diferencia del enfoque DEA.
    \item DEA es no paramétrico (toda desviación es ineficiencia); SFA es paramétrico y separa ruido aleatorio de ineficiencia.
\end{opts}

\q{En una Frontera Estocástica de Costos (SFA) para empresas de saneamiento especificada como $\ln C_i = f(Q_i, W_i; \beta) + v_i + u_i$, los dos componentes del término de error compuesto representan, respectivamente:}
\begin{opts}
    \item $v_i$ es el ruido aleatorio simétrico y $u_i \ge 0$ es la ineficiencia de costos asimétrica sobre la frontera.
    \item $v_i$ es la ineficiencia de costos estrictamente positiva y $u_i$ es el ruido estadístico simétrico con media igual a cero.
    \item Tanto $v_i$ como $u_i$ son variables aleatorias estrictamente negativas que desplazan los costos por debajo del mínimo.
    \item $v_i$ representa los costos fijos hundidos de la EPS y $u_i$ captura exclusivamente la tasa de depreciación contable.
\end{opts}

\q{En el método de corrección de sesgo de selección muestral de Heckman en dos etapas, incluir la razón inversa de Mills ($\hat{\lambda}_i$) como regresor adicional en la ecuación de resultado permite:}
\begin{opts}
    \item Corregir la correlación serial de primer orden en los errores de la ecuación de selección Probit estimada.
    \item Controlar el sesgo de selección no aleatoria y contrastarlo vía la significancia de $\hat{\lambda}_i$.
    \item Eliminar automáticamente la multicolinealidad perfecta entre los regresores continuos de la ecuación principal.
    \item Reemplazar la necesidad de contar con una variable de exclusión en la ecuación de selección sin pérdida de precisión.
\end{opts}

\subsection*{Parte II. Aplicaciones prácticas y ejercicios cuantitativos (154 a 180)}

\q{En una regresión lineal simple MCO del gasto mensual en agua ($Y$, en soles) sobre el consumo ($X$, en $m^3$), se sabe que $\bar{X} = 20$, $\bar{Y} = 85$ y la pendiente estimada es $\hat{\beta}_1 = 2.5$. ¿Cuáles son el intercepto $\hat{\beta}_0$ y el gasto predicho $\hat{Y}$ para un hogar que consume $X = 24\ m^3$?}
\begin{opts}
    \item $\hat{\beta}_0 = 30.0$ y $\hat{Y} = 90.0$
    \item $\hat{\beta}_0 = 40.0$ y $\hat{Y} = 100.0$
    \item $\hat{\beta}_0 = 35.0$ y $\hat{Y} = 95.0$
    \item $\hat{\beta}_0 = 35.0$ y $\hat{Y} = 105.0$
\end{opts}

\q{Un modelo de regresión múltiple con $n = 45$ observaciones y $k = 5$ parámetros (incluido el intercepto) reporta una suma de cuadrados totales $SCT = 500$ y una suma de cuadrados de los residuos $SCR = 100$. Calcule el $R^2$ y el $R^2$ ajustado ($\bar{R}^2$).}
\begin{opts}
    \item $R^2 = 0.80$ y $\bar{R}^2 = 0.78$
    \item $R^2 = 0.80$ y $\bar{R}^2 = 0.75$
    \item $R^2 = 0.78$ y $\bar{R}^2 = 0.76$
    \item $R^2 = 0.82$ y $\bar{R}^2 = 0.78$
\end{opts}

\q{En una regresión estimada con $n = 51$ observaciones y $5$ variables explicativas (más el intercepto, es decir $k = 6$), se obtiene un coeficiente de determinación $R^2 = 0.75$. ¿Cuál es el valor del $\bar{R}^2$ ajustado?}
\begin{opts}
    \item $0.7000$
    \item $0.7222$
    \item $0.7150$
    \item $0.7350$
\end{opts}

\q{Al estimar la elasticidad-precio de la demanda residencial de agua en un modelo log-log se obtiene $\hat{\beta}_1 = -0.64$ con error estándar $ee(\hat{\beta}_1) = 0.12$. Calcule el estadístico $t$ para contrastar la hipótesis nula de elasticidad unitaria $H_0: \beta_1 = -1$.}
\begin{opts}
    \item $t = -5.33$
    \item $t = -2.25$
    \item $t = 3.00$
    \item $t = 2.50$
\end{opts}

\q{En una muestra de $n = 66$ observaciones, un modelo no restringido con $k = 6$ parámetros alcanza $R^2_{NR} = 0.70$. Al imponer $q = 3$ restricciones lineales de exclusión, el modelo restringido obtiene $R^2_R = 0.58$. ¿Cuál es el valor del estadístico $F$?}
\begin{opts}
    \item $F = 4.00$
    \item $F = 6.00$
    \item $F = 12.00$
    \item $F = 8.00$
\end{opts}

\q{Al probar la exclusión individual ($q = 1$) de una variable cuyo coeficiente MCO estimado es $\hat{\beta}_j = 1.20$, el contraste $F$ arroja $F = 16.00$. ¿Cuál es el error estándar $ee(\hat{\beta}_j)$ de dicho coeficiente?}
\begin{opts}
    \item $0.075$
    \item $0.300$
    \item $0.240$
    \item $0.400$
\end{opts}

\q{En una ecuación de consumo de agua $\ln(Q_i) = \beta_0 + \beta_1 \ln(P_i) - 0.20\, M_i + u_i$, donde $M_i = 1$ si el hogar tiene micromedidor y $0$ en caso contrario, el efecto porcentual exacto de instalar un micromedidor sobre el consumo (use $e^{-0.20} \approx 0.8187$) es:}
\begin{opts}
    \item Una reducción exacta del $20.00\%$.
    \item Una reducción exacta del $16.67\%$.
    \item Una reducción exacta del $18.13\%$.
    \item Una reducción exacta del $22.14\%$.
\end{opts}

\q{Un estudio econométrico sobre eficiencia operativa estima la función $\hat{Y} = 50 + 12X - 0.30X^2$, donde $X$ es la antigüedad de la planta en años. ¿A qué valor de $X$ alcanza $\hat{Y}$ su nivel máximo?}
\begin{opts}
    \item $X^* = 15$ años
    \item $X^* = 40$ años
    \item $X^* = 25$ años
    \item $X^* = 20$ años
\end{opts}

\q{Se estima la regresión $\hat{Y} = 10 + 0.80X + 5.0D - 0.25(X \cdot D)$, donde $D = 1$ identifica a las EPS grandes y $D = 0$ a las pequeñas. ¿Cuál es el efecto marginal de aumentar $X$ en una unidad para una EPS grande ($D = 1$)?}
\begin{opts}
    \item $0.55$ unidades
    \item $0.80$ unidades
    \item $1.05$ unidades
    \item $5.55$ unidades
\end{opts}

\q{En el diagnóstico de multicolinealidad de una regresión múltiple, el Factor de Inflación de Varianza de la variable $X_3$ resulta igual a $VIF_3 = 8.0$. ¿Qué valor tiene el $R_3^2$ de la regresión auxiliar de $X_3$ sobre los demás regresores?}
\begin{opts}
    \item $0.800$
    \item $0.850$
    \item $0.875$
    \item $0.925$
\end{opts}

\q{Para contrastar autocorrelación hasta de orden $p = 2$ mediante la prueba de Breusch-Godfrey, se estima la regresión auxiliar con $100$ observaciones efectivas (incluyendo $\hat{u}_{t-1}$ y $\hat{u}_{t-2}$) y se obtiene $R^2_{aux} = 0.085$. ¿Cuál es el valor del estadístico $LM$ y su distribución asintótica bajo $H_0$?}
\begin{opts}
    \item $LM = 4.25$, distribuido como $\chi^2_2$
    \item $LM = 17.00$, distribuido como $\chi^2_2$
    \item $LM = 8.50$, distribuido como $\chi^2_1$
    \item $LM = 8.50$, distribuido como $\chi^2_2$
\end{opts}

\q{En una serie temporal trimestral de recaudación tarifaria se obtiene un estadístico de Durbin-Watson igual a $DW = 1.30$. Usando la aproximación estándar $DW \approx 2(1 - \hat{\rho})$, ¿cuál es el coeficiente de autocorrelación serial de primer orden estimado $\hat{\rho}$?}
\begin{opts}
    \item $\hat{\rho} = 0.35$
    \item $\hat{\rho} = 0.30$
    \item $\hat{\rho} = 0.20$
    \item $\hat{\rho} = 0.65$
\end{opts}

\q{En el modelo verdadero $Y_i = \beta_0 + 2.0 X_i^* + u_i$, la variable $X_i^*$ tiene varianza $\sigma_{X^*}^2 = 16$, pero se observa con un error de medición clásico $X_i = X_i^* + e_i$ cuya varianza es $\sigma_e^2 = 4$. ¿A qué valor converge en probabilidad el estimador MCO $\hat{\beta}_1$?}
\begin{opts}
    \item $1.20$
    \item $1.60$
    \item $1.50$
    \item $2.50$
\end{opts}

\q{Para estimar el efecto causal de un programa de capacitación ($X \in \{0,1\}$) sobre la productividad ($Y$) usando una lotería aleatoria como instrumento binario ($Z \in \{0,1\}$), se observa $\bar{Y}_{Z=1} = 42$, $\bar{Y}_{Z=0} = 36$, $\bar{X}_{Z=1} = 0.75$ y $\bar{X}_{Z=0} = 0.25$. Calcule el estimador de Wald ($\hat{\beta}_{VI}$).}
\begin{opts}
    \item $06.0$
    \item $08.0$
    \item $16.0$
    \item $12.0$
\end{opts}

\q{El modelo verdadero es $Y = \beta_0 + \beta_1 X_1 + 4.0 X_2 + u$, pero se estima por MCO omitiendo $X_2$. Si $\text{Cov}(X_1, X_2) = -3.0$ y $\text{Var}(X_1) = 10.0$, ¿cuál es la magnitud y signo del sesgo de variable omitida en $\hat{\beta}_1$?}
\begin{opts}
    \item $-1.20$ (sesgo a la baja)
    \item $+1.20$ (sesgo al alza)
    \item $-0.75$ (sesgo a la baja)
    \item $-1.50$ (sesgo a la baja)
\end{opts}

\q{Para evaluar el impacto de la micromedición en el consumo mensual de agua ($m^3$), el grupo tratado pasó de $22.0\ m^3$ antes a $17.5\ m^3$ después de la intervención, mientras que el grupo de control pasó de $20.0\ m^3$ a $18.5\ m^3$. ¿Cuál es el estimador de Diferencias en Diferencias ($\hat{\delta}_{DiD}$)?}
\begin{opts}
    \item $-4.5\ m^3$
    \item $-3.0\ m^3$
    \item $-1.5\ m^3$
    \item $-1.0\ m^3$
\end{opts}

\q{En un modelo Logit de acceso a saneamiento, el índice lineal estimado para un hogar es $X_i'\hat{\beta} = 1.0$. Sabiendo que $e^{-1} \approx 0.3679$, ¿cuál es la probabilidad predicha $\hat{P}(Y_i = 1 | X_i)$?}
\begin{opts}
    \item $0.5000$
    \item $0.6321$
    \item $0.7311$
    \item $0.8413$
\end{opts}

\q{En un modelo Logit, el coeficiente estimado de la variable $X_1$ es $\hat{\beta}_1 = 0.40$. Dado que $e^{0.40} \approx 1.4918$, un incremento de una unidad en $X_1$ (ceteris paribus) provoca que la razón de probabilidades (*odds*) del evento:}
\begin{opts}
    \item Aumente aproximadamente en $49.18\%$.
    \item Aumente aproximadamente en $40.00\%$.
    \item Aumente aproximadamente en $14.92\%$.
    \item Se multiplique exactamente por $0.40$.
\end{opts}

\q{Un modelo Probit estimado por Máxima Verosimilitud reporta una log-verosimilitud restringida (solo intercepto) de $\ln L_0 = -200$ y una log-verosimilitud no restringida de $\ln L_{ur} = -120$. Calcule el pseudo-$R^2$ de McFadden.}
\begin{opts}
    \item $0.30$
    \item $0.40$
    \item $0.60$
    \item $0.80$
\end{opts}

\q{Considere el proceso estacionario $AR(1)$ dado por $y_t = 6 + 0.6 y_{t-1} + \varepsilon_t$, donde $\varepsilon_t$ es ruido blanco con varianza $\sigma_\varepsilon^2 = 16$. ¿Cuáles son la media incondicional $E[y_t]$ y la varianza incondicional $\text{Var}(y_t)$ del proceso?}
\begin{opts}
    \item $E[y_t] = 10$ y $\text{Var}(y_t) = 20$
    \item $E[y_t] = 15$ y $\text{Var}(y_t) = 40$
    \item $E[y_t] = 15$ y $\text{Var}(y_t) = 25$
    \item $E[y_t] = 12$ y $\text{Var}(y_t) = 25$
\end{opts}

\q{Si una serie de caudales estacionaria sigue un proceso $AR(1)$ puro con coeficiente autorregresivo $\phi = 0.8$, ¿cuál es el valor teórico de su coeficiente de autocorrelación de tercer orden ($\rho_3$)?}
\begin{opts}
    \item $0.240$
    \item $0.800$
    \item $0.640$
    \item $0.512$
\end{opts}

\q{Sea el proceso de medias móviles $MA(1)$ definido por $y_t = 10 + \varepsilon_t + 0.5\varepsilon_{t-1}$, donde $\varepsilon_t$ es ruido blanco con $\sigma_\varepsilon^2 = 4$. Determine la varianza $\text{Var}(y_t)$ y la autocorrelación de primer orden $\rho_1$.}
\begin{opts}
    \item $\text{Var}(y_t) = 5.0$ y $\rho_1 = 0.50$
    \item $\text{Var}(y_t) = 5.0$ y $\rho_1 = 0.40$
    \item $\text{Var}(y_t) = 6.0$ y $\rho_1 = 0.40$
    \item $\text{Var}(y_t) = 4.5$ y $\rho_1 = 0.25$
\end{opts}

\q{En el Modelo de Corrección de Errores $\Delta y_t = 0.40\Delta x_t - 0.25(y_{t-1} - 2.0 - 1.50x_{t-1}) + \varepsilon_t$, si en el periodo $t-1$ la variable $y_{t-1}$ se ubicó $8$ unidades por encima de su equilibrio de largo plazo y $\Delta x_t = 0$, ¿cuál es el ajuste esperado $E[\Delta y_t]$ en el periodo $t$?}
\begin{opts}
    \item $-3.20$ unidades
    \item $-1.25$ unidades
    \item $-2.00$ unidades
    \item $-0.25$ unidades
\end{opts}

\q{Para un modelo estimado con $n = 100$ observaciones y $k = 5$ parámetros, se obtiene $\ln(SCR/n) = 2.00$. Usando las expresiones $AIC = \ln(SCR/n) + \frac{2k}{n}$ y $BIC = \ln(SCR/n) + \frac{k\ln(n)}{n}$ (con $\ln(100) \approx 4.605$), calcule los valores de AIC y BIC.}
\begin{opts}
    \item $AIC = 2.050$ y $BIC = 2.150$
    \item $AIC = 2.100$ y $BIC = 2.461$
    \item $AIC = 2.230$ y $BIC = 2.100$
    \item $AIC = 2.100$ y $BIC = 2.230$
\end{opts}

\q{Los residuos de una regresión con $n = 120$ observaciones presentan un coeficiente de asimetría $S = 0.5$ y una curtosis $K = 4.0$. Calcule el estadístico de normalidad de Jarque-Bera ($JB = \frac{n}{6}\left[S^2 + \frac{(K-3)^2}{4}\right]$).}
\begin{opts}
    \item $JB = 10.00$
    \item $JB = 07.50$
    \item $JB = 05.00$
    \item $JB = 15.00$
\end{opts}

\q{Se estima un modelo lineal-logarítmico del consumo mensual de agua (en $m^3$) en función del ingreso familiar: $\hat{Q}_i = 12 + 30\ln(Ingreso_i)$. Si el ingreso de un hogar aumenta en $10\%$, ¿cuál es el incremento aproximado esperado en el consumo de agua?}
\begin{opts}
    \item $0.30\ m^3$
    \item $30.0\ m^3$
    \item $3.00\ m^3$
    \item $3.00\%$
\end{opts}

\q{En un panel de $4$ años, una EPS registra un índice de cobertura de micromedición $x_{it}$ igual a $12, 18, 21$ y $25$ en los años $t = 1, 2, 3, 4$, respectivamente. Al aplicar la transformación *within* de Efectos Fijos ($\tilde{x}_{it} = x_{it} - \bar{x}_i$), ¿cuál es el valor transformado $\tilde{x}_{i3}$ correspondiente al tercer año?}
\begin{opts}
    \item $\tilde{x}_{i3} = +1.0$
    \item $\tilde{x}_{i3} = +6.0$
    \item $\tilde{x}_{i3} = +3.0$
    \item $\tilde{x}_{i3} = +2.0$
\end{opts}

\newpage
\section*{Tabla de claves rápidas}

\noindent Use esta tabla para una verificación inmediata. El solucionario con la justificación analítica de cada respuesta se presenta en la sección siguiente.

\vspace{0.5cm}
\begin{center}
\renewcommand{\arraystretch}{1.35}
\begin{tabular}{|c|c|c|c|c|c|c|c|c|c|c|}
\hline
\textbf{N.} & \textbf{1} & \textbf{2} & \textbf{3} & \textbf{4} & \textbf{5} & \textbf{6} & \textbf{7} & \textbf{8} & \textbf{9} & \textbf{10} \\
\hline
\textbf{Clave} & B & D & C & C & D & C & C & D & A & C \\
\hline
\textbf{N.} & \textbf{11} & \textbf{12} & \textbf{13} & \textbf{14} & \textbf{15} & \textbf{16} & \textbf{17} & \textbf{18} & \textbf{19} & \textbf{20} \\
\hline
\textbf{Clave} & C & B & A & B & D & C & B & A & C & B \\
\hline
\textbf{N.} & \textbf{21} & \textbf{22} & \textbf{23} & \textbf{24} & \textbf{25} & \textbf{26} & \textbf{27} & \textbf{28} & \textbf{29} & \textbf{30} \\
\hline
\textbf{Clave} & B & D & A & D & A & A & C & A & B & C \\
\hline
\textbf{N.} & \textbf{31} & \textbf{32} & \textbf{33} & \textbf{34} & \textbf{35} & \textbf{36} & \textbf{37} & \textbf{38} & \textbf{39} & \textbf{40} \\
\hline
\textbf{Clave} & A & D & C & A & A & C & B & A & A & B \\
\hline
\textbf{N.} & \textbf{41} & \textbf{42} & \textbf{43} & \textbf{44} & \textbf{45} & \textbf{46} & \textbf{47} & \textbf{48} & \textbf{49} & \textbf{50} \\
\hline
\textbf{Clave} & A & B & D & C & D & D & A & B & C & B \\
\hline
\textbf{N.} & \textbf{51} & \textbf{52} & \textbf{53} & \textbf{54} & \textbf{55} & \textbf{56} & \textbf{57} & \textbf{58} & \textbf{59} & \textbf{60} \\
\hline
\textbf{Clave} & A & A & C & A & D & C & B & C & A & C \\
\hline
\textbf{N.} & \textbf{61} & \textbf{62} & \textbf{63} & \textbf{64} & \textbf{65} & \textbf{66} & \textbf{67} & \textbf{68} & \textbf{69} & \textbf{70} \\
\hline
\textbf{Clave} & C & B & D & B & D & B & B & A & B & B \\
\hline
\textbf{N.} & \textbf{71} & \textbf{72} & \textbf{73} & \textbf{74} & \textbf{75} & \textbf{76} & \textbf{77} & \textbf{78} & \textbf{79} & \textbf{80} \\
\hline
\textbf{Clave} & D & D & A & B & B & C & D & A & B & D \\
\hline
\textbf{N.} & \textbf{81} & \textbf{82} & \textbf{83} & \textbf{84} & \textbf{85} & \textbf{86} & \textbf{87} & \textbf{88} & \textbf{89} & \textbf{90} \\
\hline
\textbf{Clave} & A & D & C & A & D & D & C & B & D & D \\
\hline
\end{tabular}
\end{center}

\vspace{0.4cm}

\newpage
\section*{Solucionario completo}

\noindent Cada respuesta incluye la clave oficial y su justificación analítica o matemática detallada.

\subsection*{Bloque A. Fundamentos teóricos y conceptuales}

\sol{1}{B}{\textbf{Tema:} Estimador matricial MCO. \textbf{Explicación:} La minimización de la suma de residuos al cuadrado $e'e = (Y - X\beta)'(Y - X\beta)$ conduce al sistema de ecuaciones normales $(X'X)\hat{\beta} = X'Y$. Si $X'X$ es invertible, la solución única es $\hat{\beta} = (X'X)^{-1} X'Y$.}

\sol{2}{D}{\textbf{Tema:} Insesgadez bajo exogeneidad. \textbf{Explicación:} Como $\hat{\beta} = \beta + (X'X)^{-1} X'u$, tomando esperanza condicional en $X$ se tiene $E[\hat{\beta}|X] = \beta + (X'X)^{-1} X' E[u|X]$. Si $E[u|X] = 0$, entonces $E[\hat{\beta}|X] = \beta$, garantizando la insesgadez.}

\sol{3}{C}{\textbf{Tema:} Rango completo y multicolinealidad perfecta. \textbf{Explicación:} Si existe una relación lineal exacta entre las variables explicativas (multicolinealidad perfecta), las columnas de $X$ son linealmente dependientes, el determinante de $X'X$ es cero y la matriz no se puede invertir.}

\sol{4}{C}{\textbf{Tema:} Supuesto de perturbaciones esféricas. \textbf{Explicación:} Perturbaciones esféricas implica que la matriz de varianza-covarianza de los errores es escalar: $\Sigma = \sigma^2 I_n$. Los elementos de la diagonal principal son constantes e idénticos (homocedasticidad) y los elementos fuera de la diagonal son nulos (no autocorrelación).}

\sol{5}{D}{\textbf{Tema:} Teorema de Gauss-Markov (BLUE). \textbf{Explicación:} Gauss-Markov demuestra que para cualquier combinación lineal $c'\beta$, el estimador lineal insesgado MCO $c'\hat{\beta}$ minimiza la varianza de estimación frente a cualquier otro estimador lineal insesgado rival.}

\sol{6}{C}{\textbf{Tema:} Estimador insesgado de la varianza residual. \textbf{Explicación:} La suma de residuos al cuadrado $e'e$ pierde $k$ grados de libertad por las $k$ condiciones ortogonales de primer orden impuestas al estimar los parámetros. Por tanto, dividiendo entre $n - k$ se obtiene $E[s^2] = \sigma^2$.}

\sol{7}{C}{\textbf{Tema:} Propiedad del R-cuadrado. \textbf{Explicación:} Al agregar regresores adicionales, el espacio de columnas de $X$ se expande y la suma de cuadrados residuales ($SCR$) nunca puede aumentar en la optimización por mínimos cuadrados. Por ende, $R^2 = 1 - SCR/STC$ nunca se reduce.}

\sol{8}{D}{\textbf{Tema:} R-cuadrado ajustado. \textbf{Explicación:} El $R^2$ ajustado corrige dividiendo $SCR$ y $STC$ entre sus respectivos grados de libertad ($n - k$ y $n - 1$). Puede incluso tomar valores negativos si el modelo es sumamente deficiente y el estadístico $F$ es menor que 1.}

\sol{9}{A}{\textbf{Tema:} Distribución normal de los coeficientes MCO. \textbf{Explicación:} Dado que $\hat{\beta}$ es una combinación lineal de las perturbaciones normales $Y$ (o $u$), hereda la distribución normal multivariada con media $\beta$ y matriz de varianzas $\sigma^2 (X'X)^{-1}$.}

\sol{10}{C}{\textbf{Tema:} Prueba t de significancia individual. \textbf{Explicación:} El cociente entre el coeficiente estimado y su error estándar estimado sigue exactamente una distribución $t$ con $n - k$ grados de libertad bajo el supuesto de normalidad de los errores, o normal asintótica en muestras grandes por el TLC.}

\sol{11}{C}{\textbf{Tema:} Prueba F conjunta de la regresión. \textbf{Explicación:} La prueba $F$ global mide si la variación explicada por los $k - 1$ regresores (excluyendo el intercepto) es significativamente mayor que la variación residual no explicada ponderada por grados de libertad.}

\sol{12}{B}{\textbf{Tema:} Interpretación de modelo log-log. \textbf{Explicación:} Dado que $\frac{d \ln(Y)}{d \ln(X)} = \frac{dY / Y}{dX / X} = \beta_1$, el parámetro mide la elasticidad directamente y es invariante a las unidades de medida de las variables.}

\sol{13}{A}{\textbf{Tema:} Interpretación de modelo log-lineal. \textbf{Explicación:} $\frac{d \ln(Y)}{dX} = \frac{1}{Y} \frac{dY}{dX} = \beta_1 \implies \frac{\Delta Y}{Y} \approx \beta_1 \Delta X$. Por tanto, multiplicando por 100 se obtiene el porcentaje de cambio aproximado en $Y$ por cada unidad de incremento en $X$.}

\sol{14}{B}{\textbf{Tema:} Consecuencias de la multicolinealidad imperfecta. \textbf{Explicación:} La multicolinealidad severa infla las varianzas muestrales de los coeficientes afectados ($VIF_j = \frac{1}{1 - R_j^2}$ muy alto). Los intervalos de confianza se ensanchan y los estadísticos $t$ se desploman, aunque el $R^2$ global del modelo sea muy elevado.}

\sol{15}{D}{\textbf{Tema:} Definición de heterocedasticidad. \textbf{Explicación:} La homocedasticidad exige $Var(u_i|X) = \sigma^2$ constante para todo $i$. Si la dispersión de las perturbaciones varía entre individuos, estratos o tamaños de empresa, existe heterocedasticidad: $Var(u_i|X) = \sigma_i^2$.}

\sol{16}{C}{\textbf{Tema:} Consecuencias de la heterocedasticidad. \textbf{Explicación:} La insesgadez no depende del supuesto de varianza constante, sino de $E[u|X]=0$. Sin embargo, la fórmula de varianza usual $\sigma^2 (X'X)^{-1}$ queda mal calculada, arrojando errores estándar engañosos y pruebas de hipótesis inválidas.}

\sol{17}{B}{\textbf{Tema:} Prueba de White. \textbf{Explicación:} Halbert White (1980) propuso una prueba que no requiere supuestos a priori sobre la forma funcional de la heterocedasticidad: se regresiona $e_i^2$ sobre variables lineales, cuadráticas e interacciones. El estadístico $LM = n R_{aux}^2$ sigue una $\chi^2$ con grados de libertad iguales al número de regresores en la ecuación auxiliar.}

\sol{18}{A}{\textbf{Tema:} Prueba de Breusch-Pagan. \textbf{Explicación:} La prueba de Breusch-Pagan evalúa si la varianza del error depende de variables observadas $Z$. La hipótesis nula es la ausencia de heterocedasticidad (homocedasticidad), de modo que los coeficientes de las variables auxiliares son nulos.}

\sol{19}{C}{\textbf{Tema:} Mínimos Cuadrados Ponderados (WLS). \textbf{Explicación:} Al transformar el modelo dividiendo cada observación entre $\sqrt{h_i}$, la perturbación transformada $u_i^* = u_i / \sqrt{h_i}$ tiene varianza $Var(u_i^*) = Var(u_i) / h_i = \sigma^2 h_i / h_i = \sigma^2$, restaurando la homocedasticidad y la eficiencia BLUE.}

\sol{20}{B}{\textbf{Tema:} Errores estándar robustos de White. \textbf{Explicación:} El estimador de Huber-White no altera $\hat{\beta}$, pero reemplaza la fórmula ingenua de varianza por la estructura sándwich consistente asintóticamente, garantizando que los estadísticos $t$ y $F$ sean fiables en muestras grandes aun con heterocedasticidad desconocida.}

\sol{21}{B}{\textbf{Tema:} Definición de proceso AR(1). \textbf{Explicación:} En un proceso autorregresivo de primer orden $AR(1)$, el error del periodo actual depende linealmente del error del periodo inmediato anterior ponderado por $\rho$ más una innovación contemporánea independiente $\varepsilon_t$.}

\sol{22}{D}{\textbf{Tema:} Consecuencias de la autocorrelación serial. \textbf{Explicación:} La correlación positiva sucesiva entre errores hace que la muestra contenga menos información independiente de la aparente ($n$ efectivo menor). MCO subestima la dispersión real, arrojando errores estándar diminutos y $t$ gigantescos que inducen a rechazar falsamente hipótesis nulas.}

\sol{23}{A}{\textbf{Tema:} Valor central de Durbin-Watson. \textbf{Explicación:} Si $\hat{\rho} = 0$, $d \approx 2(1 - 0) = 2$. Valores de $d$ cercanos a 0 indican fuerte autocorrelación positiva ($\hat{\rho} \approx 1$); valores cercanos a 4 indican fuerte autocorrelación negativa ($\hat{\rho} \approx -1$).}

\sol{24}{D}{\textbf{Tema:} Limitación de Durbin-Watson con rezagos. \textbf{Explicación:} Con variables dependientes rezagadas ($Y_{t-1}$), el estadístico $d$ está fuertemente sesgado hacia 2 (hacia no rechazar la no autocorrelación), dando una falsa impresión de limpieza estocástica. Se debe usar el estadístico $h$ o la prueba LM de Breusch-Godfrey.}

\sol{25}{A}{\textbf{Tema:} Ventajas de Breusch-Godfrey. \textbf{Explicación:} La prueba LM de Breusch-Godfrey es general: regresa los residuos $e_t$ sobre las variables explicativas y $p$ rezagos de los propios residuos ($e_{t-1}, \dots, e_{t-p}$). El estadístico $(n - p) R^2 \sim \chi^2_p$ es asintóticamente robusto frente a regresores endógenos rezagados.}

\sol{26}{A}{\textbf{Tema:} Estimador HAC de Newey-West. \textbf{Explicación:} Whitney Newey y Kenneth West diseñaron una matriz de varianzas-covarianzas que estima de manera consistente las correlaciones intertemporales no nulas y la heterocedasticidad aplicando un núcleo de ponderación decreciente (núcleo de Bartlett), garantizando que la matriz sea semidefinida positiva.}

\sol{27}{C}{\textbf{Tema:} Causas de la endogeneidad. \textbf{Explicación:} La endogeneidad rompe la exogeneidad estricta. Ocurre por: 1) omitir una variable relevante correlacionada con los regresores incluidos, 2) medir las $X$ con ruido aleatorio (sesgo de atenuación), o 3) determinación simultánea de $Y$ y $X$ en un sistema de equilibrio.}

\sol{28}{A}{\textbf{Tema:} Inconsistencia bajo endogeneidad. \textbf{Explicación:} Bajo endogeneidad, $plim \, \hat{\beta} = \beta + \frac{Cov(X, u)}{Var(X)} \neq \beta$. Al haber correlación entre $X$ y $u$, el sesgo no desaparece al aumentar el tamaño de muestra, invalidando toda interpretación causal del coeficiente.}

\sol{29}{B}{\textbf{Tema:} Condiciones de un instrumento válido. \textbf{Explicación:} La condición de relevancia exige que el instrumento mueva a la variable endógena ($Cov(Z, X) \neq 0$). La restricción de exclusión (exogeneidad) exige que $Z$ sea ortogonal al término de error poblacional ($Cov(Z, u) = 0$), incidiendo en $Y$ exclusivamente por su efecto sobre $X$.}

\sol{30}{C}{\textbf{Tema:} Procedimiento de MC2E / 2SLS. \textbf{Explicación:} La primera etapa extrae la porción exógena de la variable endógena proyectándola sobre los instrumentos ($\hat{X} = Z(Z'Z)^{-1}Z'X$). La segunda etapa utiliza únicamente este componente exógeno purificado $\hat{X}$ para estimar la ecuación estructural libre de sesgo.}

\subsection*{Bloque B. Aplicaciones cuantitativas y contrastes econométricos}

\sol{31}{A}{\textbf{Tema:} Cálculo de la pendiente MCO simple. \textbf{Explicación:} $\hat{\beta}_1 = \frac{\sum (X_i - \bar{X})(Y_i - \bar{Y})}{\sum (X_i - \bar{X})^2} = \frac{600}{400} = 1.50$.}

\sol{32}{D}{\textbf{Tema:} Cálculo de correlación y R-cuadrado. \textbf{Explicación:} $r = \frac{Cov(X, Y)}{S_X S_Y} = \frac{600}{\sqrt{400 \times 900}} = \frac{600}{\sqrt{360000}} = \frac{600}{600} = 1.00$... corrección: con 400 y 900 es $\sqrt{360000} = 600$. Entonces $r = 1.00$ y $R^2 = 1.00$.}

\sol{33}{C}{\textbf{Tema:} Cálculo de varianza residual. \textbf{Explicación:} Grados de libertad: $n - k = 45 - 5 = 40$. Varianza residual: $s^2 = \frac{SCR}{n - k} = \frac{800}{40} = 20$. Error estándar: $s = \sqrt{20} \approx 4.472$.}

\sol{34}{A}{\textbf{Tema:} Cálculo de estadístico t. \textbf{Explicación:} $t = \frac{\hat{\beta}_1}{EE(\hat{\beta}_1)} = \frac{-0.45}{0.15} = -3.00$. Como $|t| = 3.00 > 1.98$, se rechaza la hipótesis nula con un nivel de confianza superior al $99\%$.}

\sol{35}{A}{\textbf{Tema:} Cálculo del estadístico F global. \textbf{Explicación:} $F = \frac{R^2 / (k - 1)}{(1 - R^2) / (n - k)} = \frac{0.60 / 3}{(1 - 0.60) / (64 - 4)} = \frac{0.20}{0.40 / 60} = \frac{0.20}{0.00667} = 30.0$.}

\sol{36}{C}{\textbf{Tema:} Prueba F de restricciones lineales. \textbf{Explicación:} $F = \frac{(SCR_R - SCR_{NR}) / q}{SCR_{NR} / (n - k)} = \frac{(560 - 500) / 2}{500 / (106 - 6)} = \frac{60 / 2}{500 / 100} = \frac{30}{5} = 6.0$.}

\sol{37}{B}{\textbf{Tema:} Interpretación de elasticidad en demanda log-log. \textbf{Explicación:} En el modelo log-log, el coeficiente $-0.35$ es la elasticidad precio directa: $\% \Delta Q \approx -0.35 \times 10\% = -3.5\%$. Dado que $|-0.35| < 1$, la demanda de agua potable es inelástica respecto al precio.}

\sol{38}{A}{\textbf{Tema:} Cálculo de semielasticidad log-lineal. \textbf{Explicación:} $\Delta \ln(OPEX) \approx \beta_1 \Delta X = 0.04 \times 5 = 0.20$. Multiplicando por 100, los costos operativos aumentan aproximadamente un $20\%$.}

\sol{39}{A}{\textbf{Tema:} Interpretación de modelo lineal-log. \textbf{Explicación:} En especificación lineal-log, $\Delta Y \approx (\beta_1 / 100) \times (\% \Delta X) = (200 / 100) \times 1 = 2$ unidades absolutas.}

\sol{40}{B}{\textbf{Tema:} Cálculo del Factor de Inflación de Varianza (VIF). \textbf{Explicación:} $VIF_j = \frac{1}{1 - R_j^2} = \frac{1}{1 - 0.80} = \frac{1}{0.20} = 5.0$.}

\sol{41}{A}{\textbf{Tema:} Cálculo de error estándar matricial. \textbf{Explicación:} La matriz de covarianzas es $s^2 (X'X)^{-1}$. La varianza de $\hat{\beta}_2$ es $Var(\hat{\beta}_2) = s^2 \times [(X'X)^{-1}]_{33} = 25 \times 0.09 = 2.25$. El error estándar es $EE(\hat{\beta}_2) = \sqrt{2.25} = 1.50$.}

\sol{42}{B}{\textbf{Tema:} Estadístico de prueba de White. \textbf{Explicación:} $LM = n \times R_{aux}^2 = 100 \times 0.15 = 15.0$. Sigue asintóticamente una distribución $\chi^2$ con grados de libertad iguales al número de regresores en la ecuación auxiliar ($p = 5$). Como el valor crítico al $5\%$ de $\chi^2_5$ es $11.07$ y $15.0 > 11.07$, se rechaza la homocedasticidad.}

\sol{43}{D}{\textbf{Tema:} Cálculo de Breusch-Pagan. \textbf{Explicación:} $LM = 80 \times 0.10 = 8.0$. Bajo $H_0$ de homocedasticidad sigue una $\chi^2_3$. El valor crítico al $5\%$ es $7.81$. Como $8.0 > 7.81$, se rechaza la homocedasticidad.}

\sol{44}{C}{\textbf{Tema:} Transformación analítica de WLS. \textbf{Explicación:} La desviación estándar del error es $\sqrt{\sigma^2 X_i^2} = \sigma X_i$. Dividiendo toda la ecuación entre $X_i$: $\frac{Y_i}{X_i} = \beta_0 (1/X_i) + \beta_1 + u_i^*$, donde $Var(u_i^*) = Var(u_i/X_i) = \sigma^2$, logrando homocedasticidad perfecta.}

\sol{45}{D}{\textbf{Tema:} Cálculo de Durbin-Watson a partir de rho. \textbf{Explicación:} $d \approx 2(1 - \hat{\rho}) = 2(1 - 0.60) = 2(0.40) = 0.80$.}

\sol{46}{D}{\textbf{Tema:} Regla de decisión de Durbin-Watson. \textbf{Explicación:} Si el estadístico calculado $d$ es estrictamente menor que el límite inferior tabulado $d_L$ ($0.80 < 1.63$), se rechaza categóricamente la hipótesis nula $H_0: \rho = 0$ a favor de la alternativa de autocorrelación serial positiva.}

\sol{47}{A}{\textbf{Tema:} Estimación de rho para Cochrane-Orcutt. \textbf{Explicación:} Despejando de $d \approx 2(1 - \hat{\rho}) \implies 1 - \hat{\rho} = d / 2 \implies \hat{\rho} = 1 - d/2 = 1 - 1.30 / 2 = 1 - 0.65 = 0.35$.}

\sol{48}{B}{\textbf{Tema:} Transformación de Prais-Winsten en t=1. \textbf{Explicación:} A diferencia de Cochrane-Orcutt (que descarta la primera observación), Prais-Winsten conserva la primera observación multiplicándola por $\sqrt{1 - \rho^2}$, preservando la eficiencia en muestras pequeñas.}

\sol{49}{C}{\textbf{Tema:} Cálculo de sesgo de variable omitida. \textbf{Explicación:} La fórmula del sesgo de variable omitida es $E[\tilde{\beta}_1] = \beta_1 + \beta_2 \frac{Cov(X_1, X_2)}{Var(X_1)} = \beta_1 + \beta_2 \tilde{\gamma}_1 = 4 + 2(0.5) = 4 + 1.0 = 5.0$.}

\sol{50}{B}{\textbf{Tema:} Cálculo del estimador de Variables Instrumentales. \textbf{Explicación:} El estimador de variables instrumentales simple es el cociente de covarianzas muestrales: $\hat{\beta}_{IV} = \frac{Cov(Z, Y)}{Cov(Z, X)} = \frac{150}{75} = 2.00$.}

\sol{51}{A}{\textbf{Tema:} Regla de instrumentos débiles de Stock-Yogo. \textbf{Explicación:} La regla empírica estándar de Stock y Yogo establece que si el estadístico $F$ de los instrumentos excluidos en la primera etapa es menor a 10 ($4.2 < 10$), los instrumentos son débiles, generando sesgo significativo y distorsión de tamaño en los contrastes de Wald.}

\sol{52}{A}{\textbf{Tema:} Mecánica de la prueba de Hausman en regresión. \textbf{Explicación:} Bajo $H_0$ de exogeneidad, el residuo de primera etapa $\hat{v}$ no debe tener poder explicativo sobre $Y$ una vez controlado $X_1$. Si su coeficiente es significativo en la regresión aumentada, $X_1$ y el error estructural están correlacionados (endogeneidad confirmada).}

\sol{53}{C}{\textbf{Tema:} Prueba de sobreidentificación de Sargan. \textbf{Explicación:} La hipótesis nula de la prueba de Sargan / Hansen es que todos los instrumentos son exógenos y válidos. Como el estadístico calculado $1.20 < 5.99$, no se rechaza $H_0$, respaldando la validez del instrumental utilizado.}

\sol{54}{A}{\textbf{Tema:} Trampa de las variables ficticias. \textbf{Explicación:} La suma de las 4 variables dummy es exactamente igual al vector columna de unos de la constante ($D_1 + D_2 + D_3 + D_4 = 1$), provocando dependencia lineal exacta. Para evitarlo se debe omitir una categoría base o suprimir la constante.}

\sol{55}{D}{\textbf{Tema:} Efecto marginal con variable de interacción. \textbf{Explicación:} La derivada parcial respecto a la educación es $\frac{\partial \ln(W)}{\partial Educ} = \beta_1 + \beta_3 Mujer$. Para hombres ($Mujer=0$) es $\beta_1$; para mujeres ($Mujer=1$) es $\beta_1 + \beta_3$.}

\sol{56}{C}{\textbf{Tema:} Propiedad de regresión estandarizada. \textbf{Explicación:} Al tener todas las variables media cero, la ecuación evaluada en las medias exige $\bar{Y} = \hat{\beta}_0 + \sum \hat{\beta}_j \bar{X}_j \implies 0 = \hat{\beta}_0 + 0 \implies \hat{\beta}_0 = 0$.}

\sol{57}{B}{\textbf{Tema:} Propiedades de la Hat Matrix. \textbf{Explicación:} $H$ proyecta el vector $Y$ sobre el subespacio generado por las columnas de $X$ ($\hat{Y} = HY$). Como operador de proyección ortogonal, es simétrica e idempotente, y su traza es igual al rango de $X$, que es $k$.}

\sol{58}{C}{\textbf{Tema:} Ortogonalidad de residuos MCO. \textbf{Explicación:} Las condiciones de primer orden de MCO exigen $\frac{\partial (e'e)}{\partial \beta} = -2X'e = 0 \implies X'e = 0$. Si hay intercepto, una columna de $X$ es de unos, garantizando $\sum e_i = 0$ y covarianza muestral nula con los regresores.}

\sol{59}{A}{\textbf{Tema:} Descomposición de varianza e intercepto. \textbf{Explicación:} El término cruzado $\sum (\hat{Y}_i - \bar{Y})e_i$ se anula únicamente si el modelo contiene constante (o si los residuos suman cero). Sin constante, $STC \neq SEC + SCR$ y el $R^2$ estándar carece de sentido matemático.}

\sol{60}{C}{\textbf{Tema:} Cálculo de intervalo de confianza. \textbf{Explicación:} $IC_{95\%} = \hat{\beta}_1 \pm t_{crit} \times EE(\hat{\beta}_1) = 2.50 \pm 2.00 \times 0.25 = 2.50 \pm 0.50 = [2.00, \; 3.00]$.}

\sol{61}{C}{\textbf{Tema:} Varianza del error de predicción individual. \textbf{Explicación:} La predicción individual incorpora dos fuentes de incertidumbre: la estimación de los parámetros $\beta$ (varianza $\sigma^2 x_0'(X'X)^{-1}x_0$) más la varianza de la nueva perturbación $u_0$ que afectará al individuo ($\\sigma^2$), sumando $\sigma^2 [1 + x_0'(X'X)^{-1}x_0]$.}

\sol{62}{B}{\textbf{Tema:} Comparación entre AIC y BIC. \textbf{Explicación:} La penalización del AIC es $2k$, mientras que la del BIC es $k \ln(n)$. Como para $n \ge 8$ se cumple $\ln(n) > 2$, el BIC castiga más drásticamente la sobreparametrización y es un criterio consistente para identificar el modelo verdadero.}

\sol{63}{D}{\textbf{Tema:} Prueba de normalidad de Jarque-Bera. \textbf{Explicación:} Bajo la hipótesis nula de normalidad, $S \approx 0$ y $K \approx 3$. El estadístico de Jarque-Bera combina las desviaciones cuadráticas de asimetría y exceso de curtosis, siguiendo asintóticamente una distribución $\chi^2$ con 2 grados de libertad.}

\sol{64}{B}{\textbf{Tema:} Prueba RESET de Ramsey. \textbf{Explicación:} Si la forma funcional lineal es correcta, las potencias de los valores predichos ($\hat{Y}^2, \hat{Y}^3$) no deberían tener poder explicativo residual. Un estadístico $F$ significativo indica no linealidad o sesgo por variables omitidas.}

\sol{65}{D}{\textbf{Tema:} Propiedad de MCO restringido. \textbf{Explicación:} Imponer restricciones reduce el espacio de búsqueda en la minimización de residuos. Por ende, la suma de residuos al cuadrado restringida nunca puede ser menor que la no restringida: $SCR_R \ge SCR_{NR}$ (igualdad solo si la muestra satisface exactamente la restricción).}

\subsection*{Bloque C. Casos con implicancias regulatorias y sector saneamiento}

\sol{66}{B}{\textbf{Tema:} Modelación de demanda de agua en SUNASS. \textbf{Explicación:} La forma funcional log-log permite interpretar inmediatamente los coeficientes estimados como las elasticidades directas de consumo frente a cambios tarifarios y variaciones en el ingreso de los hogares, parámetros esenciales para diseñar la estructura tarifaria quinquenal.}

\sol{67}{B}{\textbf{Tema:} Endogeneidad en tarifas por bloques. \textbf{Explicación:} Bajo esquemas de bloques crecientes, la tarifa marginal no es exógena: depende del volumen consumido. Si un choque no observado eleva el consumo ($u > 0$), el usuario salta de bloque tarifario ($P$ sube), generando correlación positiva espuria entre precio y error.}

\sol{68}{A}{\textbf{Tema:} Instrumentación en demanda de agua. \textbf{Explicación:} La tarifa del primer bloque básico es fijada exógenamente por el regulador y es independiente del consumo discrecional final del hogar, funcionando como un instrumento válido que correlaciona con la tarifa promedio sin compartir choques de demanda.}

\sol{69}{B}{\textbf{Tema:} Frontera estocástica en regulación de EPS. \textbf{Explicación:} SFA especifica el residuo como $\varepsilon_i = v_i + u_i$, donde $v_i \sim N(0, \sigma_v^2)$ es ruido estocástico blanco (clima, geografía) y $u_i \ge 0$ es ineficiencia técnica pura. Esto evita atribuir a mala gestión costos derivados de choques exógenos adversos.}

\sol{70}{B}{\textbf{Tema:} Economías de densidad en saneamiento. \textbf{Explicación:} Las economías de densidad de clientes/producto miden cómo responde el costo al expandir el volumen suministrado dentro de una red de distribución ya existente (sin extender kilómetros de tubería), siendo típicamente crecientes en servicios de saneamiento.}

\sol{71}{D}{\textbf{Tema:} Heterocedasticidad por tamaño de empresa. \textbf{Explicación:} En muestras con disparidad extrema de tamaño (Sedapal atiende a más de 10 millones de habitantes mientras que pequeñas EPS atienden a 20 mil), la variabilidad de los errores es función directa del tamaño de la empresa, requiriendo estimación ponderada o errores robustos.}

\sol{72}{D}{\textbf{Tema:} Errores robustos HC3 en muestras pequeñas. \textbf{Explicación:} Davidson y MacKinnon demostraron que en muestras con menos de 100 observaciones, el estimador HC0 de White subestima la verdadera varianza en puntos de alto leverage. El estimador HC3 divide los residuos al cuadrado entre $(1 - h_{ii})^2$, corrigiendo el sesgo muestral.}

\sol{73}{A}{\textbf{Tema:} Modelación SARIMA en hidrología urbana. \textbf{Explicación:} Las series hidrológicas presentan alta memoria temporal y patrones estacionales marcados (estiaje y avenidas). Los modelos SARIMA $(p,d,q)(P,D,Q)_{12}$ capturan la dinámica de corto plazo y la estacionalidad anual del recurso hídrico.}

\sol{74}{B}{\textbf{Tema:} Newey-West en series de tiempo climáticas. \textbf{Explicación:} Los choques climáticos como El Niño generan correlaciones residuales que persisten a lo largo de múltiples trimestres. Los errores HAC de Newey-West blindan la inferencia estadística evitando rechazos espurios de significancia causal.}

\sol{75}{B}{\textbf{Tema:} Sesgo de atrición en paneles de medidores. \textbf{Explicación:} Si la salida de medidores del panel no es un proceso puramente aleatorio (ej. se rompen más rápido en casas de consumo masivo o con mayor presión de agua), la muestra superviviente se vuelve no representativa, sesgando los coeficientes de demanda.}

\sol{76}{C}{\textbf{Tema:} Diferencias en Diferencias en saneamiento. \textbf{Explicación:} Diff-in-Diff controla por factores macroeconómicos y climáticos temporales comunes que afectan a toda la ciudad: calcula la diferencia antes-después en el grupo tratado y le resta la diferencia antes-después del grupo de control: $\hat{\delta}_{DD} = (\bar{Y}_{T,post} - \bar{Y}_{T,pre}) - (\bar{Y}_{C,post} - \bar{Y}_{C,pre})$.}

\sol{77}{D}{\textbf{Tema:} Supuesto de tendencias paralelas. \textbf{Explicación:} Si en el periodo previo al tratamiento el grupo de control y tratamiento ya exhibían trayectorias divergentes, la diferencia observada post-tratamiento no puede atribuirse con rigor causal a la instalación del medidor, invalidando la evaluación.}

\sol{78}{A}{\textbf{Tema:} Regresión Discontinua (RDD) en SISFOH. \textbf{Explicación:} Alrededor del punto de corte $c_0$, los hogares justo por encima y justo por debajo del umbral tienen características socioeconómicas prácticamente idénticas (asignación cuasialeatoria), permitiendo estimar el efecto causal del subsidio tarifario de forma muy robusta.}

\sol{79}{B}{\textbf{Tema:} Limitaciones del Modelo Lineal de Probabilidad. \textbf{Explicación:} El MLP ($P(Y=1|X) = X\beta$) genera predicciones absurdas ($P < 0$ o $P > 1$) en los extremos de la distribución y su error tiene varianza heterocedástica $Var(u_i|X) = X\beta(1 - X\beta)$. Logit y Probit confinan la probabilidad entre 0 y 1 mediante funciones de distribución acumuladas.}

\sol{80}{D}{\textbf{Tema:} Multicolinealidad en infraestructura de EPS. \textbf{Explicación:} Las EPS con más kilómetros de red tienen inevitablemente más conexiones domiciliarias y producen más volumen de agua. Esta colinealidad infla las varianzas de los coeficientes individuales, aunque el modelo conjunto prediga con gran exactitud el costo total.}

\sol{81}{A}{\textbf{Tema:} Causalidad inversa en ANF y costos. \textbf{Explicación:} Existe retroalimentación simultánea entre el gasto en mantenimiento de redes y las fugas físicas de agua: más fugas demandan más reparaciones, pero menos mantenimiento provoca más roturas de tuberías, requiriendo variables instrumentales para identificar el efecto causal.}

\sol{82}{D}{\textbf{Tema:} Corrección de selección de Heckman. \textbf{Explicación:} Los hogares con conexión domiciliaria no son una muestra aleatoria de la población urbana (autoselección). Estimar MCO ignorando a los no conectados introduce sesgo de selección. El método de Heckman estima primero la probabilidad de tener conexión (Probit) e inserta la Razón Inversa de Mills en la ecuación de gasto.}

\sol{83}{C}{\textbf{Tema:} Ventaja de Efectos Fijos en datos de panel de EPS. \textbf{Explicación:} Cada EPS opera bajo condiciones geográficas y topográficas únicas que no varían con el tiempo. El estimador Within de Efectos Fijos elimina $\alpha_i$ diferenciando respecto a la media individual, evitando sesgos por variables omitidas territoriales.}

\sol{84}{A}{\textbf{Tema:} Prueba de Hausman en datos de panel. \textbf{Explicación:} Bajo la hipótesis nula, ambos estimadores son consistentes pero RE es más eficiente. Si $H_0$ se rechaza, la heterogeneidad no observada $\alpha_i$ está correlacionada con los regresores, haciendo que RE sea sesgado e inconsistente, por lo que la teoría exige adoptar Efectos Fijos.}

\sol{85}{D}{\textbf{Tema:} Estimación de la Beta CAPM en WACC regulatorio. \textbf{Explicación:} La Beta mide el riesgo sistemático no diversificable del negocio de saneamiento. Se estima regresionando los retornos bursátiles de empresas del sector sobre el índice de mercado (S\&P 500) y ajustando por apalancamiento financiero.}

\sol{86}{D}{\textbf{Tema:} Cálculo de Beta de Hamada en regulación. \textbf{Explicación:} $\beta_L = 0.50 \times [1 + (1 - 0.30) \times 0.60] = 0.50 \times [1 + 0.70 \times 0.60] = 0.50 \times [1 + 0.42] = 0.50 \times 1.42 = 0.71$.}

\sol{87}{C}{\textbf{Tema:} Modelos de conteo en roturas de tuberías. \textbf{Explicación:} Las fallas de red son eventos discretos no negativos. El modelo de Poisson modela conteos pero asume equi-dispersión ($Var = Media$). Si los datos presentan sobredispersión por incidentes en racimo, el modelo Binomial Negativo introduce un parámetro de dispersión gamma para inferencia correcta.}

\sol{88}{B}{\textbf{Tema:} Detección econométrica de cárteles en licitaciones. \textbf{Explicación:} Los cárteles coordinan ofertas de resguardo que se agrupan en torno a márgenes prefijados, generando una distribución de posturas con dispersión artificialmente baja y patrones de simetría incompatibles con la competencia no cooperativa.}

\sol{89}{D}{\textbf{Tema:} Econometría de Nowcasting en saneamiento. \textbf{Explicación:} El Nowcasting (predicción del presente inmediato) utiliza modelos de factores dinámicos o MIDAS para explotar indicadores de alta frecuencia disponibles de inmediato, permitiendo anticipar caídas en la cobranza de las EPS antes de que se publiquen los estados financieros oficiales.}

\sol{90}{D}{\textbf{Tema:} Poder estadístico y errores Tipo I y II en calidad de agua. \textbf{Explicación:} El diseño del muestreo bacteriológico es una decisión econométrica/estadística crucial: rechazar erróneamente el cumplimiento impone costos judiciales injustos a la EPS (Tipo I), mientras que no detectar contaminación expone a la población a epidemias diarreicas (Tipo II), requiriendo un poder estadístico ($1 - \beta$) riguroso.}

\subsection*{Bloque D. Ampliación: Teoría y Aplicaciones Prácticas (91 a 180)}

\sol{91}{A}{La matriz $M = I_n - P_X$ proyecta sobre el complemento ortogonal de $X$: es simétrica ($M'=M$), idempotente ($M^2=M$), aniquila a $X$ ($MX=0$) y $\text{tr}(M)=n-k$.}
\sol{92}{B}{El teorema FWL demuestra que $\hat{\beta}_1 = (X_1' M_2 X_1)^{-1} X_1' M_2 Y$, es decir, MCO sobre las variables previamente depuradas de la influencia lineal de $X_2$.}
\sol{93}{C}{Sin intercepto en $X$, la condición de primer orden $X'\hat{u}=0$ no garantiza $\sum \hat{u}_i = 0$, por lo que $SCT \neq SCE + SCR$ y $1 - SCR/SCT$ puede ser negativo.}
\sol{94}{D}{Por la ley de expectativas iteradas, $E[u|X]=0 \Rightarrow E[X'u]=0$, pero no al revés. $E[u|X]=0$ da insesgadez exacta; $E[X'u]=0$ garantiza consistencia asintótica.}
\sol{95}{B}{Bajo normalidad, MCO coincide con MV y alcanza la cota de Cramér-Rao, siendo el Mejor Estimador Insesgado (MIE/BUE) incluso frente a estimadores no lineales.}
\sol{96}{C}{Si $X_2^* = c X_2$, entonces $\hat{\beta}_2^* = \hat{\beta}_2/c$ y $ee(\hat{\beta}_2^*) = ee(\hat{\beta}_2)/c$; por tanto, $t_2$, $\hat{Y}$ y $R^2$ permanecen exactamente invariantes.}
\sol{97}{D}{$u_i = Y_i - X_i'\beta$ es inobservable y esférico ($\sigma^2 I$), en tanto $\hat{u} = Mu$ es observable y tiene matriz de varianzas $\sigma^2 M$ (ni independiente ni homocedástica).}
\sol{98}{A}{Puede demostrarse algebraicamente que añadir un regresor reduce la varianza residual estimada $s^2 = SCR/(n-k)$ (y por ende sube $\bar{R}^2$) si y solo si $|t_{k+1}| > 1$.}
\sol{99}{C}{Al pasar de $D=0$ a $D=1$, $\frac{Y_1 - Y_0}{Y_0} = e^{\delta} - 1$. Así, el cambio porcentual exacto es $100(e^{\delta}-1)\%$; la regla $100\delta\%$ es solo una aproximación de primer orden.}
\sol{100}{D}{Derivando respecto a $X$: $\frac{\partial E[Y|X]}{\partial X} = \beta_1 + 2\beta_2 X$. Igualando a cero se obtiene el punto crítico $X^* = -\beta_1/(2\beta_2)$, que es máximo pues $2\beta_2 < 0$.}
\sol{101}{A}{El cuadrado de una variable $t_{n-k}$ sigue exactamente una distribución $F_{1,n-k}$. Por ello, para $q=1$ restricción lineal, el contraste $F$ y el contraste $t$ bilateral son idénticos.}
\sol{102}{B}{Como $ee(\hat{\beta}_j) \propto 1/\sqrt{n}$, en muestras muy grandes cualquier efecto real distinto de cero —por diminuto que sea— genera $|t|$ muy grande y $p$-valor cercano a $0$.}
\sol{103}{D}{La prueba de Chow contrasta si tanto el intercepto como todas las pendientes permanecen constantes entre submuestras, lo cual equivale a probar conjuntamente la dummy de periodo y todas sus interacciones.}
\sol{104}{A}{En el modelo lineal normal con restricciones lineales se cumple exactamente $W \ge LR \ge LM$, y asintóticamente los tres estadísticos son equivalentes bajo $H_0$ con distribución $\chi^2_q$.}
\sol{105}{B}{Sobre-especificar el modelo añadiendo una variable con $\beta=0$ no introduce sesgo ni inconsistencia en MCO, pero reduce eficiencia (infla las varianzas) si correlaciona con los demás regresores.}
\sol{106}{C}{Con $\text{Cov}(X_1, X_2)=0$, el sesgo de variable omitida en $\hat{\beta}_1$ es cero; sin embargo, $\beta_2 X_2$ pasa al término de error, inflando $s^2$ y sesgando al alza los errores estándar.}
\sol{107}{A}{Ambos criterios se minimizan. Como para $n \ge 8$ se cumple $\ln(n) > 2$, BIC penaliza más la complejidad que AIC, favorece modelos más parsimoniosos y es consistente en selección.}
\sol{108}{B}{Con $n$ fijo, reducir $\alpha$ estrecha la región de rechazo, dificultando rechazar $H_0$: baja el Error Tipo I, pero sube el Error Tipo II ($\beta$) y disminuye la potencia ($1-\beta$).}
\sol{109}{C}{La multicolinealidad infla las varianzas individuales (bajando los $|t|$), pero no afecta la capacidad explicativa conjunta ($R^2$ y $F$ global altos) ni viola el carácter MELI de MCO.}
\sol{110}{D}{$R_j^2$ es el $R^2$ de regresar $X_j$ contra las demás explicativas; $1/(1-R_j^2)$ es el $VIF_j$, que mide cuánto se infla la varianza de $\hat{\beta}_j$ por colinealidad.}
\sol{111}{B}{La prueba de White detecta divergencias entre la matriz de varianzas clásica y la robusta, lo cual ocurre tanto ante heterocedasticidad pura como ante mala especificación funcional.}
\sol{112}{C}{Goldfeld-Quandt ordena según $X_j$, elimina $c$ observaciones intermedias para acentuar la diferencia de varianzas, estima MCO en ambos extremos y contrasta $F = SCR_2/SCR_1$.}
\sol{113}{D}{Al sustituir $\Omega$ por una estimación $\hat{\Omega}$, FGLS pierde la garantía exacta de insesgadez en muestras finitas, pero conserva consistencia y eficiencia asintótica sobre MCO.}
\sol{114}{A}{Los procesos ARCH/GARCH modelan agrupamientos de volatilidad (*volatility clustering*): la varianza condicional $\sigma_t^2$ varía en el tiempo aunque la varianza incondicional sea finita y constante.}
\sol{115}{C}{Como $DW \approx 2(1 - \hat{\rho})$, cuando $\hat{\rho} < 0$ se tiene $DW > 2$, acercándose a $4$ cuando $\hat{\rho} \to -1$.}
\sol{116}{D}{Con $Y_{t-1}$ en el lado derecho, DW se sesga hacia 2; se emplea el estadístico $h$ de Durbin (para $AR(1)$) o preferentemente Breusch-Godfrey (válido para cualquier orden $p$).}
\sol{117}{A}{Cochrane-Orcutt descarta $t=1$ al cuasi-diferenciar ($y_t - \hat{\rho}y_{t-1}$). Prais-Winsten preserva $t=1$ escalando por $\sqrt{1-\hat{\rho}^2}$, lo que mejora la eficiencia en muestras pequeñas.}
\sol{118}{B}{Al igual que los errores robustos de White, Newey-West no modifica $\hat{\beta}_{MCO}$; solo reemplaza la matriz de varianzas por un estimador consistente ante heterocedasticidad y autocorrelación.}
\sol{119}{D}{Un error clásico en $Y$ se absorbe en $u$, manteniendo insesgadez con mayor varianza. Un error clásico en $X$ correlaciona el regresor observado con el error compuesto, sesgando $\hat{\beta}_1$ hacia cero.}
\sol{120}{A}{La condición de orden (contar exclusiones: $K - k \ge m - 1$) es necesaria mas no suficiente; la condición de rango sobre la matriz de parámetros reducidos es necesaria y suficiente.}
\sol{121}{B}{Imbens y Angrist (1994) demostraron que con efectos heterogéneos, validez del instrumento y monotonicidad, el estimador VI identifica el LATE: el efecto causal promedio para los *compliers*.}
\sol{122}{C}{Con instrumentos débiles ($\text{Cov}(Z,X) \approx 0$), MC2E tiene una distribución no normal, varianza enorme y un sesgo en muestras finitas que lo arrastra hacia el estimador sesgado de MCO.}
\sol{123}{A}{Con $L = m$ (exactamente identificado), los residuos VI cumplen $Z'\hat{u}_{VI} = 0$ por definición algebraica ($0$ grados de sobreidentificación), por lo que la exclusión debe defenderse teóricamente.}
\sol{124}{B}{MC2E es un caso particular de GMM que usa la matriz de pesos óptima solo bajo homocedasticidad. Con heterocedasticidad y sobreidentificación, GMM en dos etapas es asintóticamente más eficiente.}
\sol{125}{C}{La prueba de Sargan asume errores homocedásticos; bajo heterocedasticidad o errores agrupados, el estadístico $J$ de Hansen (basado en GMM) mantiene distribución asintótica $\chi^2_{L-m}$.}
\sol{126}{D}{Controlar por variables determinadas después del tratamiento (*bad controls*) introduce sesgo de selección al condicionar en un resultado o colisionador, distorsionando el efecto causal total.}
\sol{127}{B}{La transformación intragrupo resta la media temporal de cada individuo. Si $z_{it} = z_i$ para todo $t$, entonces $z_{it} - \bar{z}_i = 0$ y queda colinealmente absorbida por el efecto fijo $\alpha_i$.}
\sol{128}{C}{Por el Teorema de Frisch-Waugh-Lovell, proyectar fuera las dummies individuales produce exactamente las desviaciones respecto a la media individual; por ello $\hat{\beta}_{within} \equiv \hat{\beta}_{LSDV}$.}
\sol{129}{D}{Si $u_{it}$ es ruido blanco, FE es más eficiente que FD (pues $\Delta u_{it}$ tiene autocorrelación $-0.5$). Si $u_{it}$ es caminata aleatoria, $\Delta u_{it} = e_{it}$ es serialmente no correlacionado y FD es más eficiente.}
\sol{130}{A}{Como $\theta = 1 - \sqrt{\frac{\sigma_u^2}{\sigma_u^2 + T\sigma_{\alpha}^2}}$, si $T\sigma_{\alpha}^2 \gg \sigma_u^2$ entonces $\theta \to 1$ y RE se aproxima a FE; si $\sigma_{\alpha}^2 \to 0$, $\theta \to 0$ y RE tiende a MCO agrupado.}
\sol{131}{C}{Agrupar (*cluster*) los errores estándar a nivel de unidad (EPS) garantiza inferencia asintóticamente válida ante heterocedasticidad y correlación serial arbitraria dentro de cada clúster.}
\sol{132}{D}{En paneles cortos con $y_{it-1}$, la media $\bar{u}_i$ contiene $u_{it-1}$, correlacionando $\tilde{y}_{it-1}$ con $\tilde{u}_{it}$ (sesgo de Nickell). Arellano-Bond usa primeras diferencias e instrumentos en niveles rezagados.}
\sol{133}{A}{FE es robusto a la atrición si la selección depende de regresores $X_{it}$ o de $\alpha_i$ (ya que $\alpha_i$ se elimina), pero se sesga si la salida depende de los choques idiosincráticos $u_{it}$.}
\sol{134}{B}{En un proceso TS ($y_t = \alpha + \delta t + u_t$ con $u_t$ estacionario), los shocks mueren y $y_t$ vuelve a la tendencia. En un proceso DS ($y_t = \delta + y_{t-1} + \varepsilon_t$), cada shock se acumula permanentemente.}
\sol{135}{D}{Un proceso $AR(p)$ es estacionario si todas las raíces de $\phi(z)=1-\phi_1 z - \dots - \phi_p z^p = 0$ satisfacen $|z| > 1$ (en el caso $AR(1)$, $1-\phi_1 z=0 \Rightarrow |z|=1/|\phi_1|>1 \Leftrightarrow |\phi_1|<1$).}
\sol{136}{A}{En un proceso $MA(q)$, la memoria se corta en el rezago $q$ (FAC $= 0$ para $k > q$) y la FACP decae gradualmente; en un $AR(p)$ ocurre exactamente el patrón simétrico inverso.}
\sol{137}{B}{La prueba DF simple asume errores no correlacionados. La versión Aumentada (ADF) añade rezagos de $\Delta y_t$ para limpiar la autocorrelación residual, manteniendo valores críticos de Dickey-Fuller.}
\sol{138}{C}{El Teorema de Representación de Granger establece que dos series $I(1)$ están cointegradas si y solo si admiten un MCE, implicando además causalidad de Granger en al menos una dirección.}
\sol{139}{A}{Si $y_{t-1}$ excede su valor de equilibrio ($u_{t-1} > 0$), la corrección requiere que $\Delta y_t$ baje; por ello se exige $\alpha < 0$ (típicamente $-1 < \alpha < 0$ para convergencia monótona).}
\sol{140}{B}{La causalidad de Granger evalúa precedencia predictiva (contrastada mediante prueba $F$/Wald sobre los rezagos de $X$ en la ecuación de $Y$), no causalidad estructural filosófica.}
\sol{141}{C}{Cholesky impone una estructura recursiva triangular inferior sobre los impactos contemporáneos; si la matriz de covarianzas residuales no es diagonal, cambiar el orden de las variables altera las IRF.}
\sol{142}{D}{Como $\frac{\partial P(Y=1|X)}{\partial X_j} = g(X'\beta)\beta_j$ y $g(\cdot)>0$, el signo coincide con $\beta_j$, pero la magnitud varía según $X$; se resume mediante el efecto en la media (MEM) o el promedio (AME).}
\sol{143}{B}{En Logit, $\ln\left(\frac{p}{1-p}\right) = X'\beta$, de modo que aumentar $X_j$ en una unidad multiplica los *odds* $\frac{p}{1-p}$ exactamente por $e^{\beta_j}$.}
\sol{144}{C}{Con censura (solución de esquina o límite de medición) conocemos $X_i$ de toda la muestra aunque $Y_i$ se apile en el límite; con truncamiento se excluyen de la base los individuos fuera del rango.}
\sol{145}{D}{Poisson impone equidispersión ($\text{Var}=E$). Bajo sobredispersión, Poisson cuasi-MV sigue siendo consistente para $\beta$ si la media está bien especificada, pero los $ee$ estándar son demasiado pequeños; Binomial Negativa modela explícitamente la sobredispersión.}
\sol{146}{A}{Como en modelos discretos $0 > \ln L_{ur} \ge \ln L_0$, el cociente $\ln L_{ur}/\ln L_0 \in (0, 1]$ y $R^2_{McF} = 1 - \ln L_{ur}/\ln L_0 \in [0, 1)$, evaluando mejora relativa en log-verosimilitud.}
\sol{147}{C}{PSM asume que, condicionando en $X$ (o en $p(X)$), el tratamiento es independiente de los resultados potenciales (CIA) y que existe traslape de probabilidades $0 < P(D=1|X) < 1$ (soporte común).}
\sol{148}{D}{En *Fuzzy RDD*, el indicador de cruzar el umbral actúa como instrumento válido para el tratamiento efectivo cerca del corte: $\hat{\tau}_{FRD} = \frac{\lim_{x\downarrow c}E[Y|x] - \lim_{x\uparrow c}E[Y|x]}{\lim_{x\downarrow c}E[D|x] - \lim_{x\uparrow c}E[D|x]}$.}
\sol{149}{A}{La validez de RDD exige que los individuos no puedan manipular con precisión su puntaje alrededor del corte; McCrary contrasta la continuidad de la densidad de la *running variable* en el umbral.}
\sol{150}{B}{Goodman-Bacon (2021) y Callaway-Sant'Anna muestran que con adopción escalonada y efectos dinámicos heterogéneos, TWFE realiza comparaciones 'prohibidas' usando ya tratados como controles, pudiendo asignar pesos negativos.}
\sol{151}{D}{DEA construye una frontera determinística por programación lineal (sensible a outliers/ruido); SFA estima una frontera paramétrica descomponiendo la desviación en ruido simétrico y término de ineficiencia.}
\sol{152}{A}{En SFA de costos (minimización), $v_i \sim N(0, \sigma_v^2)$ es ruido aleatorio bidireccional y $u_i \ge 0$ (seminormal, exponencial o normal truncada) mide el sobrecosto por ineficiencia.}
\sol{153}{B}{En Heckman en dos etapas, $\hat{\lambda}_i = \phi(z_i'\hat{\gamma})/\Phi(z_i'\hat{\gamma})$ corrige la media condicional truncada; si su coeficiente es significativamente distinto de cero, existe sesgo de selección.}
\sol{154}{C}{La recta MCO pasa por $(\bar{X}, \bar{Y})$: $\hat{\beta}_0 = \bar{Y} - \hat{\beta}_1\bar{X} = 85 - 2.5(20) = 35.0$. Luego, $\hat{Y} = 35.0 + 2.5(24) = 95.0$.}
\sol{155}{A}{$R^2 = 1 - \frac{100}{500} = 0.80$. Además, $\bar{R}^2 = 1 - \frac{SCR/(n-k)}{SCT/(n-1)} = 1 - \frac{100/40}{500/44} = 1 - 0.22 = 0.78$.}
\sol{156}{B}{Aplicando $\bar{R}^2 = 1 - (1 - R^2)\frac{n-1}{n-k} = 1 - (0.25)\frac{50}{45} = 1 - 0.2778 = 0.7222$.}
\sol{157}{C}{El estadístico es $t = \frac{\hat{\beta}_1 - \beta_{1,0}}{ee(\hat{\beta}_1)} = \frac{-0.64 - (-1.00)}{0.12} = \frac{0.36}{0.12} = 3.00$.}
\sol{158}{D}{Usando $F = \frac{(R^2_{NR} - R^2_R)/q}{(1 - R^2_{NR})/(n - k)} = \frac{(0.70 - 0.58)/3}{(1 - 0.70)/(66 - 6)} = \frac{0.04}{0.005} = 8.00$.}
\sol{159}{B}{Para $q = 1$ se cumple $F = t^2$, luego $|t| = \sqrt{16} = 4.0$. Así, $ee(\hat{\beta}_j) = \frac{|\hat{\beta}_j|}{|t|} = \frac{1.20}{4.0} = 0.300$.}
\sol{160}{C}{El cambio porcentual exacto es $100(e^{\hat{\delta}} - 1)\% = 100(e^{-0.20} - 1)\% = 100(0.8187 - 1)\% = -18.13\%$.}
\sol{161}{D}{Derivando respecto a $X$: $\frac{d\hat{Y}}{dX} = 12 - 0.60X = 0 \Rightarrow X^* = \frac{12}{0.60} = 20$ años (máximo porque $-0.60 < 0$).}
\sol{162}{A}{El efecto marginal es $\frac{\partial \hat{Y}}{\partial X} = 0.80 - 0.25D$. Evaluando en $D = 1$ se obtiene $0.80 - 0.25 = 0.55$ unidades.}
\sol{163}{C}{Como $VIF_3 = \frac{1}{1 - R_3^2} = 8$, se despeja $1 - R_3^2 = \frac{1}{8} = 0.125 \Rightarrow R_3^2 = 0.875$.}
\sol{164}{D}{El estadístico de Breusch-Godfrey es $LM = n \cdot R^2_{aux} = 100 \times 0.085 = 8.50$, el cual sigue asintóticamente una distribución $\chi^2_p = \chi^2_2$ bajo $H_0$.}
\sol{165}{A}{Despejando de $DW \approx 2(1 - \hat{\rho})$: $\hat{\rho} \approx 1 - \frac{DW}{2} = 1 - \frac{1.30}{2} = 1 - 0.65 = 0.35$.}
\sol{166}{B}{Bajo error de medición clásico en $X$, $\text{plim}\,\hat{\beta}_1 = \beta_1 \frac{\sigma_{X^*}^2}{\sigma_{X^*}^2 + \sigma_e^2} = 2.0 \times \frac{16}{16 + 4} = 2.0 \times 0.80 = 1.60$.}
\sol{167}{D}{El estimador de Wald es la razón entre el efecto intención de tratar y la primera etapa: $\hat{\beta}_{Wald} = \frac{\bar{Y}_1 - \bar{Y}_0}{\bar{X}_1 - \bar{X}_0} = \frac{42 - 36}{0.75 - 0.25} = \frac{6}{0.50} = 12.0$.}
\sol{168}{A}{El sesgo de variable omitida es $\text{Sesgo}(\hat{\beta}_1) = \beta_2 \frac{\text{Cov}(X_1, X_2)}{\text{Var}(X_1)} = 4.0 \times \frac{-3.0}{10.0} = -1.20$.}
\sol{169}{B}{$\hat{\delta}_{DiD} = (\bar{Y}_{T,1} - \bar{Y}_{T,0}) - (\bar{Y}_{C,1} - \bar{Y}_{C,0}) = (17.5 - 22.0) - (18.5 - 20.0) = -4.5 - (-1.5) = -3.0\ m^3$.}
\sol{170}{C}{En el modelo Logit, $\hat{P}(Y=1|X) = \frac{1}{1 + e^{-X'\hat{\beta}}} = \frac{1}{1 + e^{-1}} = \frac{1}{1.3679} \approx 0.7311$.}
\sol{171}{A}{El *odds ratio* es $e^{\hat{\beta}_1} = e^{0.40} \approx 1.4918$, lo que equivale a un incremento porcentual en los *odds* de $(1.4918 - 1) \times 100\% = 49.18\%$.}
\sol{172}{B}{El pseudo-$R^2$ de McFadden es $R^2_{McF} = 1 - \frac{\ln L_{ur}}{\ln L_0} = 1 - \frac{-120}{-200} = 1 - 0.60 = 0.40$.}
\sol{173}{C}{Para un $AR(1)$ estacionario: $E[y_t] = \frac{c}{1-\phi} = \frac{6}{0.4} = 15$, y $\text{Var}(y_t) = \frac{\sigma_\varepsilon^2}{1-\phi^2} = \frac{16}{1-0.36} = \frac{16}{0.64} = 25$.}
\sol{174}{D}{En un proceso $AR(1)$ estacionario, la autocorrelación de orden $k$ cumple $\rho_k = \phi^k$. Para $k = 3$: $\rho_3 = (0.8)^3 = 0.512$.}
\sol{175}{B}{Para un $MA(1)$: $\text{Var}(y_t) = \sigma_\varepsilon^2(1 + \theta^2) = 4(1 + 0.25) = 5.0$ y $\rho_1 = \frac{\theta}{1 + \theta^2} = \frac{0.5}{1.25} = 0.40$.}
\sol{176}{C}{Con $\Delta x_t = 0$ y desequilibrio $(y_{t-1} - y^*_{t-1}) = +8$, el ajuste esperado es $E[\Delta y_t] = -0.25(8) = -2.00$ unidades (corrige $25\%$ del desvío por periodo).}
\sol{177}{D}{$AIC = 2.00 + \frac{2(5)}{100} = 2.100$. $BIC = 2.00 + \frac{5(4.605)}{100} = 2.00 + 0.230 = 2.230$.}
\sol{178}{A}{Sustituyendo en la fórmula: $JB = \frac{120}{6}\left[(0.5)^2 + \frac{(4.0 - 3.0)^2}{4}\right] = 20(0.25 + 0.25) = 10.00$.}
\sol{179}{C}{En un modelo lin-log, $\Delta Q \approx \left(\frac{\hat{\beta}_1}{100}\right) \%\Delta X = \left(\frac{30}{100}\right)(10) = 3.00\ m^3$.}
\sol{180}{D}{La media temporal de la EPS es $\bar{x}_i = \frac{12 + 18 + 21 + 25}{4} = 19.0$. Por tanto, para $t = 3$ se obtiene $\tilde{x}_{i3} = 21 - 19 = +2.0$.}


\end{document}
